\documentclass[a4paper,11pt]{ltjsarticle}

% ========== package ========== %
% --- 数式環境 --- %
\usepackage{amsmath}
\usepackage{amssymb}
\usepackage{mathtools}
% --- 数式環境 END --- %
% ブラケット
\usepackage{braket}
% ベクトルの太文字
\usepackage{bm}
% 添え字左右上下 \fourIdx{左上}{左下}{右上}{右下}{ベース} で使う．Iは大文字
\usepackage{fouridx}
% 色
\usepackage{color}
% 余白指定
\usepackage{geometry}
% --- 索引（hyperref より前に読み込む） --- %
\usepackage{makeidx}
\makeindex
% --- ページ内リンク --- %
\usepackage{hyperref}
\hypersetup{colorlinks=true,linkcolor=blue,hypertexnames=true,plainpages=false}
% --- ページ内リンク END --- %
% 花文字
\usepackage{mathrsfs}
% 画像
\usepackage{graphicx}
% tcolorbox
\usepackage{tcolorbox}
% appendix
\usepackage{appendix}
% --- 参考文献 --- %
\usepackage{etoolbox}
\usepackage{fontspec} % fontspecは読み込んでOKですが、設定順序に注意が必要です
%% --- 参考文献フォント --- %%
  \IfFontExistsTF{Times New Roman}{
    \newfontfamily\bibfontface[Ligatures=TeX]{Times New Roman}
  }{
    \newfontfamily\bibfontface[Ligatures=TeX]{FreeSerif} % 代替フォント
  }
%% --- 参考文献フォント END --- %%
%% --- 参考文献リンク位置 --- %%
  \BeforeBeginEnvironment{thebibliography}{%
    \phantomsection
    \addcontentsline{toc}{section}{参考文献}
  }
  \AtBeginEnvironment{thebibliography}{%
    \normalfont\bibfontface
  }
%% --- 参考文献のリンク位置 --- %%
% --- 参考文献 END --- %
% --- 索引の見出し・目次リンク --- %
\renewcommand{\indexname}{\hyperlink{toc}{索引}}
\BeforeBeginEnvironment{theindex}{%
  \clearpage
  \phantomsection
  \addcontentsline{toc}{section}{索引}
}
% 項目は英語キーでアルファベット順に整列し、各ページ番号を本文へリンクする。
% page.<表示番号> の宛先を重複させない（通し番号、または Roman 前付＋Arabic 本文）。
% hypertexnames=false だけで重複を回避すると索引が誤ページへ飛ぶため使用しない。
% 用語の出現箇所に記す例：
% \index{decoherence@Decoherence(デコヒーレンス)|hyperpage}
% 生成手順：LuaLaTeX → BibTeX（文献がある場合）→ makeindex → LuaLaTeX 2回。
% --- 索引 END --- %
% ========== package END ========== %

% ========== newcommand ========== %
% --- セクションから目次へのリンク --- %
\NewDocumentCommand{\sectionlink}{s m}{%
  \IfBooleanTF{#1}{%
    \section*{\hyperlink{toc}{#2}}%
    \phantomsection\addcontentsline{toc}{section}{#2}%
  }{\section[#2]{\hyperlink{toc}{#2}}}%
}
\newcommand{\subsectionlink}[1]{\subsection[#1]{\hyperlink{toc}{#1}}}
% セクション番号リンク
\makeatletter
\renewcommand{\@seccntformat}[1]{%
  \protect\hyperlink{toc}{\csname the#1\endcsname}\quad
}
\makeatother

\numberwithin{equation}{section}
% ==========- newcommand END ========== %
% Explicit original equation tags leave amsmath's counter unchanged.
% Step it separately to keep the automatic PDF destinations unique.
\BeforeBeginEnvironment{equation}{\stepcounter{equation}}
\allowdisplaybreaks
%===========================================================
% --- タイトル --- %
%===========================================================
\geometry{left=25mm, right=25mm, top=30mm, bottom=30mm}
% 日本語題名 → 英語原題 → arXiv abs → INSPIRE-HEP literature の順。
% URL が空の行は表示しない。対象文献がない場合は出版社・DOI 等で代用する。
\newcommand{\papertitlejp}{量子 Brownian 運動への経路積分アプローチ}
\newcommand{\papertitleen}{Path integral approach to quantum Brownian motion}
\newcommand{\paperarxivurl}{https://doi.org/10.1016/0378-4371(83)90013-4} % 例：https://arxiv.org/abs/2508.14262v2
\newcommand{\paperarxivlabel}{Physica A 121 (1983) 587--616 (DOI)} % 例：arXiv:2508.14262v2 [quant-ph]
\newcommand{\paperinspireurl}{https://inspirehep.net/literature/176411} % 例：https://inspirehep.net/literature/文献番号
\newcommand{\paperinspirelabel}{INSPIRE-HEP: 176411} % 例：INSPIRE-HEP、または代用先の出版社・DOI
\title{%
  \bfseries\papertitlejp\\[0.5em]
  {\normalfont\large\papertitleen}%
  \ifdefempty{\paperarxivurl}{}{\\[0.4em]{\normalfont\normalsize
    \href{\paperarxivurl}{\paperarxivlabel}}}%
  \ifdefempty{\paperinspireurl}{}{\\[0.4em]{\normalfont\normalsize
    \href{\paperinspireurl}{\paperinspirelabel}}}%
}
% Original PDF page 1 (printed p. 587).
% Metadata commands are invoked by the main document.
\newcommand{\paperauthors}{%
  A.O. Caldeira\thanks{A.O. Caldeira が1980年9月に University of Sussex に D.Phil. 学位取得のため提出した学位論文の一部に基づく。}\\[0.3em]
  {\small Instituto de Física ``Gleb Wataghin'', Universidade Estadual de Campinas,}\\
  {\small 13.100 Campinas -- SP -- Brasil}\\[0.8em]
  A.J. Leggett\\[0.3em]
  {\small School of Mathematical and Physical Sciences, University of Sussex,}\\
  {\small Brighton BN1 9QH, Sussex, UK}%
}
\newcommand{\originalreceiptdates}{1982年8月25日受理\quad 1983年1月4日改訂}
\newcommand{\paperabstract}{%
Feynman と Vernon の影響汎関数法を、任意の温度における Brownian 運動の研究に適用する。対象とする系と環境との散逸的相互作用について特定の模型を選ぶことにより、影響汎関数を閉じた形で評価し、現象論的な粘性係数など少数のパラメータで表すことができる。$\hbar\to0$ の極限では、影響汎関数の形式から得られる結果が古典的な Fokker--Planck 方程式に帰着することを示す。任意の強さの減衰を受ける単純調和振動子の場合には、任意の温度で、系がある特定の種類の純粋状態から出発するときの完全な密度行列 $\rho(x,x',t)$ の時間発展を明示的に表す式を得る。得られた結果を、量子力学における散逸の問題に対する他の方法の結果と比較する。%
}

\author{\paperauthors}
\hypersetup{pdftitle={量子 Brownian 運動への経路積分アプローチ},pdfauthor={A.O. Caldeira and A.J. Leggett}}
\pdfstringdefDisableCommands{\def\hyperlink#1#2{#2}}
\date{2026年10月29日}

%===========================================================
% --- 以下本文 --- %
%===========================================================
\begin{document}
\maketitle
\begin{center}\small\originalreceiptdates\end{center}
\begin{abstract}
\paperabstract
\end{abstract}
\clearpage
\hypertarget{toc}{}
\tableofcontents
\clearpage
% Keep original reference numbering even when a citation appears later in prose.
\nocite{ref01,ref02,ref03,ref04,ref05,ref06,ref07,ref08,ref09,ref10,ref11,ref12,ref13,ref14,ref15,ref16,ref17,ref18,ref19,ref20,ref21,ref22,ref23,ref24,ref25,ref26,ref27,ref28,ref29,ref30,ref31,ref32,ref33,ref34,ref35,ref36,ref37,ref38}
% Original PDF pages 1--10 (printed pp. 587--596).
\sectionlink{序論}

自然界には、古典的な領域では次の形の運動方程式によって十分に記述できる系が数多く存在する。
\begin{equation}
 m\ddot{x}+\eta\dot{x}+v'(x)=F(t).
 \tag{1.1}\label{eq:1.1}
\end{equation}
これは、例えば粘性流体中に浸されたコロイド程度の大きさの粒子を表すことができる。この場合、$m$ は粒子の質量、$\eta$ は減衰定数、$v(x)$ は粒子に作用するポテンシャルであり、$F(t)$ はいわゆる揺動力である。この力は
\begin{equation}
 \begin{aligned}
 \langle F(t)\rangle&=0,\\
 \langle F(t)F(t')\rangle&=2\eta kT\delta(t-t')
 \end{aligned}
 \tag{1.2}\label{eq:1.2}
\end{equation}
という関係に従う。ここで $\langle\ \rangle$ は、同じように準備された系の集団についての統計平均を表す。関係 \eqref{eq:1.2} を満たす式 \eqref{eq:1.1} は、よく知られた Langevin 方程式であり、すでに半世紀以上にわたって Brownian 運動の理論の基礎として用いられてきた\cite{ref01}。

% Original PDF page 2 (printed p. 588).
他のあらゆる現象論的方程式と同様に、式 \eqref{eq:1.1} の有効範囲は限られている。系の長時間での振る舞いに関心がある場合には、Langevin 方程式を用いるのが妥当である。ここで長時間とは、対象とする系に結合した熱浴の緩和時間に比べて長い時間を意味する。先に挙げた例では、粘性媒質中の巨視的粒子は古典論によって十分に記述できるため、量子効果を気にする必要はない。しかし、他の多くの場合には事情は異なる。例えば LCR 回路を考えると、コンデンサーの極板上の電荷についての方程式は式 \eqref{eq:1.1} の形をとることが分かっている。温度が非常に低くなると、この系で量子効果が観測されると予想される。したがって、減衰を含む運動方程式と量子化の手続きをどう両立させるか、という極めて自然な問題が生じる。

この問題の起源は、標準的な量子化の手続きが、対象とする系について Hamiltonian または Lagrangian 関数のいずれかが存在することに基づいているという事実にある。一方、陽に時間に依存しないいかなる Lagrangian または Hamiltonian にも、古典的な Lagrange 方程式または Hamilton 方程式を適用して式 \eqref{eq:1.1} を得ることはできないことがよく知られている。時間に依存する関数を用いれば、標準的な量子化の手続きを直接利用することはできる\cite{ref02}。しかし、この方法では不確定性原理に関する幾つかの問題に直面せざるを得ない\cite{ref03}。過去20年から30年にわたり、多くの人々がこの問題に答えようとしてきた。用いられた方法は多様であるが、これらの試みはすべて二つの主要な範疇に分かれる。すなわち、新しい量子化の方式を探すか、あるいは系と熱浴を合わせて扱う方法を用いるかである。以下では、それぞれの範疇における比較的重要な研究の一部を挙げる。

1970年代の初めに、Kostin\cite{ref04} は非線形 Schrödinger 方程式を用いた理論を提案した。この理論は重ね合わせの原理に反するだけでなく、定常的な減衰状態のような、非常に議論の多い結果も示す。Nelson の確率量子化\cite{ref05} に基づいて、K. Yasue\cite{ref06} は Kostin の方程式を導いた。これが可能なのは、Nelson の量子化では Hamiltonian や Lagrangian の代わりに運動方程式そのものを知ればよいからである。ここでの主要な問題は、Nelson の理論が正しいかどうかということである。

Dekker\cite{ref07} は、複素変数を用いる正準量子化の手続きを開発した。密度演算子の Wigner 変換に対する Fokker--Planck 方程式など、幾つかの興味深い結果を再現するにもかかわらず、この理論には不明瞭な点がある。例えば、彼は位置 $x(t)$ の方程式に陽な雑音源を導入し、さらに運動量の方程式にも別の雑音源を導入している。そのようにする物理的な理由は見いだせない。Dekker の研究における幾つかの点には後で立ち戻る。

% Original PDF page 3 (printed p. 589).
われわれの知る限り、これまでに最も成功している方法は Senitzky\cite{ref08} の研究によって切り開かれたものである。そこで彼は、対象とする系と熱浴との相互作用を陽に考慮した。その研究では、特に共振器内の電磁場モードの減衰を扱っている。彼の方策は、熱浴の演算子を消去することによって Heisenberg 描像で実行される。この理論は後に Mori\cite{ref09} によって一般化された。同じ方向に沿って、Zwanzig\cite{ref10} と Nakajima\cite{ref11} は Schrödinger 描像で研究し、密度演算子に対する一般化されたマスター方程式を得た。この種の方法についての優れた総説は文献 \cite{ref12,ref13} にある。

この問題に対するもう一つの方法は、古典的な Langevin 方程式 \eqref{eq:1.1} を量子の場合へ一般化しようとするものである。これを行う最も単純な方法は、式 \eqref{eq:1.1} 自体は変えずに、式 \eqref{eq:1.2} によって表された古典的な力の相関を
\begin{equation}
 \langle F(t)F(t')\rangle
 =\frac{1}{2\pi}\int e^{-i\omega(t-t')}\eta\hbar\omega
 \coth\left(\frac{\hbar\omega}{2kT}\right)\,d\omega
 \tag{1.2$'$}\label{eq:1.2prime}
\end{equation}
という式で置き換えることである。ここでは $\omega$ に、例えば $\Omega$ という適切な切断を設ける。$|t-t'|\gg\Omega^{-1}$ という時間間隔に関心があるならば、式 \eqref{eq:1.2prime} は高温では、すなわち $kT\gg\hbar|t-t'|^{-1}$ では、事実上式 \eqref{eq:1.2} と等価であるが、低温ではそれと異なる。ごく最近、この方法は Koch ら\cite{ref14} によって Josephson 接合の低温での振る舞いの解析に用いられた。その理論的基礎も、最近 Benguria と Kac\cite{ref15} によって論じられている。後者の議論から明らかなように、調和振動子という特殊な場合を超えて「量子 Langevin 方程式」がどこまで有効であるかという一般的な問題は、大きく未解決のままである。

本論文で採用する一般的な方法は、文献 \cite{ref08,ref09,ref10,ref11,ref12,ref13} の研究と多くの共通点を持つ。すなわち、散逸系それ自体を量子化することは試みず、初めから複雑な環境と相互作用するものとして扱う。まさにこの相互作用が散逸を生じさせるのである。系と環境からなる完全な「宇宙」は閉じたものとして扱うことができるので、それに標準的な量子化の手続きを適用することには、もちろん何の支障もない。その後、環境の座標を消去し、散逸系だけについて閉じた方程式を得る。この消去を行うには、当然ながら系と環境との相互作用について十分に単純な模型を選ばなければならない。さらに、その模型は古典的な高温極限で古典的な Brownian 運動の結果を再現するように構成されていなければならない。

関連した問題を研究した論文は数多い。Ford ら\cite{ref16} は、二次のポテンシャルを通じて互いに相互作用する同種粒子の列を研究した。そのうちの一つに着目し、古典的な領域でそれを Brownian 粒子として振る舞わせるには、粒子間ポテンシャルがどのような形でなければならないかを調べた。Iche と Nozières\cite{ref17} は、環境中の重い粒子について断熱展開を構成したが、得られた結果は主として形式的なものである。

% Original PDF page 4 (printed p. 590).
Davies\cite{ref18} は、系の粒子の拡散的な振る舞いや、Markoff 型の Fokker--Planck 方程式などについて、幾つかの優れた厳密な結果を証明した。しかし、それらが成り立つのは、減衰が極めて弱い場合だけである。Schwinger\cite{ref19} は、熱浴に弱く結合した量子振動子の運動に自身の作用原理を適用し、振動数のずれと一定の減衰という問題を注意深く研究した。

本論文の目的は、質量や摩擦係数など、原理的には実験で測定できる現象論的パラメータだけを含む、散逸的な量子力学系の閉じた方程式を、どこまで得られるかを調べることである。このために、系と環境との相互作用について特定の模型を設定し、適切な極限で Brownian 運動の古典的方程式を再現するという条件から、模型のパラメータを定める\cite{ref20}。その上で、より一般的な場合、例えば量子的な場合に何が導かれるかを調べる。密度行列の伝播関数について、相互作用のない場合にすでに現れる式よりもわずかに複雑なだけの一般的な形式的表式を、汎関数積分の形で書き下せることを示す。また、調和振動子という特殊な場合には、任意の減衰の強さと任意の温度について、それを明示的に評価できることを示す。応用として、ある種類の純粋状態から出発する減衰調和振動子の時間発展を明示的に研究する。一貫して用いる手法は、量子力学の汎関数積分による定式化であり、より具体的には Feynman と Vernon の「影響汎関数」の手法である。

この研究の当初の動機は、ポテンシャル障壁を通り抜ける量子力学的トンネル現象に対する散逸の効果を理解したいということであった。汎関数積分法そのものはこの問題に非常によく適しているが、本論文で展開する特定の方法はそうではないことが分かった。逆に、トンネル問題に最も適した方法、すなわち「虚時間」表示を用い、4点関数ではなく2点関数を考える方法は、古典的に到達可能な領域での系の時間発展を論じるのにはあまり適していない。この理由から、文献 \cite{ref21} に示したトンネル問題の扱いとは、本論文の議論を分けるのが妥当であると考える。それでも両者に共通する問題は多く、特に散逸機構についてのわれわれの模型の一般性と現実性の程度、および速度に依存する結合によって生じる修正がそれに当たる。こうした問題の一部は文献 \cite{ref21} でさらに詳しく論じている。したがって本論文の各所で、この研究を相互参照する。

第2節では、影響汎関数の Feynman--Vernon 理論\cite{ref22,ref23} を概説する。第3節と第4節では、熱浴の具体的な模型について影響汎関数を評価する。これによって、適切な状況での Brownian 粒子の縮約密度演算子の伝播関数が得られる。第5節では、量子的な表式の $\hbar\to0$ の極限を調べることにより、Brownian 運動の古典的方程式を再現することを扱う。第3節と第4節の結果の応用を第6節に示す。特に、時刻ゼロにおいて Gaussian 型の、必ずしも最小不確定性ではない純粋状態の波束で記述される減衰調和振動子を考え、その後のすべての時刻について密度行列を明示的に計算する。最後に第7節で結論を述べる。

% Original PDF page 5 (printed p. 591).
\sectionlink{Feynman--Vernon 理論}

第2の系 B、すなわち熱浴と相互作用する系 A を考え、次の Hamiltonian で記述する。
\begin{equation}
 H=H_{\mathrm A}+H_{\mathrm I}+H_{\mathrm B}.
 \tag{2.1}\label{eq:2.1}
\end{equation}
ここで
\begin{equation}
 H_{\mathrm A}=-\frac{\hbar^2}{2M}\frac{\partial^2}{\partial x^2}+v(x)
 \tag{2.2}\label{eq:2.2}
\end{equation}
は、質量 $M$、座標 $x$ の粒子によって表される系の Hamiltonian であり、
\begin{equation}
 H_{\mathrm B}=-\frac{\hbar^2}{2m}\frac{\partial^2}{\partial\bm R^2}
 +\frac12\sum_{i\ne j}v_{\mathrm B}(R_i,R_j)
 \tag{2.3}\label{eq:2.3}
\end{equation}
は、座標 $R_i$、質量 $m$ の $N$ 個の粒子からなる熱浴の Hamiltonian である。$N$ 成分ベクトル $\bm R$ は $\bm R=(R_1,\ldots,R_N)$ を表す。最後に、$H_{\mathrm I}$ は
\begin{equation}
 H_{\mathrm I}=\sum_i v_i(x,R_i)
 \tag{2.4}\label{eq:2.4}
\end{equation}
で与えられる相互作用 Hamiltonian である。

Feynman--Vernon 理論\cite{ref22,ref23} を述べるに当たっての出発点は、時刻 $t$ における系と熱浴を合わせた密度演算子の積分形
\begin{equation}
 \rho(t)=\exp\left(-\frac{iHt}{\hbar}\right)\rho(0)
 \exp\left(\frac{iHt}{\hbar}\right)
 \tag{2.5}\label{eq:2.5}
\end{equation}
である。これは座標表示では
\begin{equation}
 \begin{aligned}
 \langle x\bm R|\rho(t)|y\bm Q\rangle
 ={}&\int dx'\,dy'\,d\bm R'\,d\bm Q'
 \left\langle x,\bm R\left|\exp\left(-\frac{iHt}{\hbar}\right)\right|x'\bm R'\right\rangle\\
 &\quad\times\langle x'\bm R'|\rho(0)|y'\bm Q'\rangle
 \left\langle y'\bm Q'\left|\exp\left(\frac{iHt}{\hbar}\right)\right|y\bm Q\right\rangle
 \end{aligned}
 \tag{2.6}\label{eq:2.6}
\end{equation}
と書くことができる。

% Original PDF page 6 (printed p. 592).
式 \eqref{eq:2.6} で $H$ を含む項は、容易に次のように認められる\cite{ref22}。
\begin{equation}
 \begin{aligned}
 \left\langle x\bm R\left|\exp\left(-\frac{iHt}{\hbar}\right)\right|x'\bm R'\right\rangle
 &\equiv K(x,\bm R,t;x',\bm R',0)\\
 &=\iint\mathcal D x\,\mathcal D\bm R\,
 \exp\left(\frac{i}{\hbar}S[x,\bm R]\right),
 \end{aligned}
 \tag{2.7}\label{eq:2.7}
\end{equation}
また、
\begin{equation}
 \begin{aligned}
 \left\langle y'\bm Q'\left|\exp\left(\frac{iHt}{\hbar}\right)\right|y\bm Q\right\rangle
 &\equiv K^*(y,\bm Q,t;y',\bm Q',0)\\
 &=\iint\mathcal D y\,\mathcal D\bm Q\,
 \exp\left(-\frac{i}{\hbar}S[y,\bm Q]\right).
 \end{aligned}
 \tag{2.8}\label{eq:2.8}
\end{equation}
すべての汎関数積分は、端点が $x(t)=x$、$x(0)=x'$、$y(t)=y$、$y(0)=y'$、$\bm R(t)=\bm R$、$\bm R(0)=\bm R'$、$\bm Q(t)=\bm Q$、$\bm Q(0)=\bm Q'$ である経路 $x(t')$、$y(t')$、$\bm R(t')$、$\bm Q(t')$ について評価する。式 \eqref{eq:2.7} と \eqref{eq:2.8} における作用 $S$ は
\begin{equation}
 S=S_{\mathrm A}+S_{\mathrm I}+S_{\mathrm B}=\int_0^t L\,dt'
 \tag{2.9}\label{eq:2.9}
\end{equation}
で与えられ、ここで
\begin{equation}
 \begin{aligned}
 L=L_{\mathrm A}+L_{\mathrm I}+L_{\mathrm B}
 ={}&\frac12 M\dot{x}^{,2}-v(x)-\sum_i v_i(x,R_i)\\
 &+\sum_i\frac12m\dot R_i^{,2}-\frac12\sum_{i\ne j}v_{\mathrm B}(R_i,R_j)
 \end{aligned}
 \tag{2.10}\label{eq:2.10}
\end{equation}
である。こうして式 \eqref{eq:2.6} は
\begin{equation}
 \begin{aligned}
 \langle x\bm R|\rho(t)|y\bm Q\rangle
 ={}&\int dx'\,dy'\,d\bm Q'\,d\bm R'
 K(x,\bm R,t;x',\bm R',0)\\
 &\times K^*(y,\bm Q,t;y',\bm Q',0)
 \langle x'\bm R'|\rho(0)|y'\bm Q'\rangle
 \end{aligned}
 \tag{2.11}\label{eq:2.11}
\end{equation}
と書くことができる。

この密度演算子の表式は、系と熱浴を合わせた全体の振る舞いを記述する。しかし、そのすべての情報を求めているわけではない。熱浴の具体的な振る舞いによらず、系 A の性質を知りたいのである。必要なのは、熱浴が系 A に及ぼす影響だけである。したがって、実際に求めている量は全密度演算子ではなく、いわゆる縮約密度演算子である\cite{ref03,ref12,ref13}。これは、式 \eqref{eq:2.11} においてすべての環境の座標についてトレースをとれば容易に得られる。すなわち、
\begin{equation}
 \begin{aligned}
 \widetilde\rho(x,y,t)
 &\equiv\int d\bm R\,\langle x\bm R|\rho(t)|y\bm R\rangle\\
 &=\int dx'\,dy'\,d\bm R'\,d\bm Q'\,d\bm R\,
 K(x,\bm R,t;x',\bm R',0)\\
 &\qquad\times K^*(y,\bm R,t;y',\bm Q',0)
 \langle x'\bm R'|\rho(0)|y'\bm Q'\rangle.
 \end{aligned}
 \tag{2.12}\label{eq:2.12}
\end{equation}
ここで、初期の全密度演算子が
\begin{equation}
 \rho(0)=\rho_{\mathrm A}(0)\rho_{\mathrm B}(0)
 \tag{2.13}\label{eq:2.13}
\end{equation}
で与えられると仮定する。

% Original PDF page 7 (printed p. 593).
ここで $\rho_{\mathrm B}(0)$ と $\rho_{\mathrm A}(0)$ は、それぞれ熱浴と系が互いに相互作用していないときの密度演算子である。$t=0$ の後では、両者とも時間とともに変化してよい。式 \eqref{eq:2.13} を式 \eqref{eq:2.12} に代入し、式 \eqref{eq:2.7} から \eqref{eq:2.10} を用いると、
\begin{equation}
 \widetilde\rho(x,y,t)=\int dx'\,dy'\,
 J(x,y,t;x',y',0)\rho_{\mathrm A}(x',y',0)
 \tag{2.14}\label{eq:2.14}
\end{equation}
を得る。ここで
\begin{equation}
 J(x,y,t;x',y',0)=\iint\mathcal D x\,\mathcal D y\,
 \exp\left(\frac{iS_{\mathrm A}[x]}{\hbar}\right)
 \exp\left(-\frac{iS_{\mathrm A}[y]}{\hbar}\right)\mathcal F[x,y]
 \tag{2.15}\label{eq:2.15}
\end{equation}
であり、
\begin{equation}
 \begin{aligned}
 \mathcal F[x,y]={}&\int d\bm R'\,d\bm Q'\,d\bm R\,
 \rho_{\mathrm B}(\bm R',\bm Q',0)\iint\mathcal D\bm R\,\mathcal D\bm Q\\
 &\times\exp\left\{\frac{i}{\hbar}\left(
 S_{\mathrm I}[x,\bm R]-S_{\mathrm I}[y,\bm Q]
 +S_{\mathrm B}[\bm R]-S_{\mathrm B}[\bm Q]\right)\right\}
 \end{aligned}
 \tag{2.16}\label{eq:2.16}
\end{equation}
が、いわゆる影響汎関数である\cite{ref22,ref23}。

対象とする系と環境との間に相互作用がなければ、影響汎関数は1に等しい。このとき式 \eqref{eq:2.15} は、一つは時間の順方向へ、他方は逆方向へ進む二つの伝播関数の積に帰着する。これは、閉じた系では密度演算子が波動関数とその複素共役との積として書けるからである。相互作用を入れると、この二つの伝播は自由には起こらなくなり、代わって影響汎関数がそれらを結び付ける。

式 \eqref{eq:2.15} は、われわれが展開する Brownian 運動の量子論における中心的な公式である。

\sectionlink{模型}

原理的には問題は解けている。式 \eqref{eq:2.3} と \eqref{eq:2.4} の粒子間ポテンシャル $v_i(x,R_i)$ と $v_{\mathrm B}(R_i,R_j)$、および $\rho_{\mathrm B}(0)$ で記述される熱浴の初期状態が分かれば、式 \eqref{eq:2.16} の影響汎関数を求めることができる。しかし、経験から分かるように、現実はそれほど単純ではない。多体系の分野で厳密に解ける問題はごく少ない。問題を完全に一般的なまま厳密に解けると考えるのは望みがない。必要なのは、構成上解くことのできる特定の模型を用いることであり、それから、多くの他の系にも拡張できると考えられる幾つかの一般的な結論を引き出すことを期待する。

熱浴は、絶対最小値をとる点 $\bm R_0=(R_{10},R_{20},\ldots,R_{N0})$ を持つ一般的なポテンシャル $v(R_1,\ldots,R_N)$ を通じて相互作用する粒子からなると仮定する。ここで、一般的な Hamiltonian \eqref{eq:2.2} で記述される系 A を熱浴に結合し、系に対する熱浴の線形応答だけを考えればよいという意味で、その結合は弱いものとする。

% Original PDF page 8 (printed p. 594).
このとき熱浴は系によってごく弱く摂動されるだけなので、熱浴を調和近似で記述できる\footnote{散逸系の一般的な記述としてこの模型を正当化できるかという問題は、文献 \cite{ref21} でさらに論じている。}。二つの系の相互作用は仮定により線形であり、系は一般的なポテンシャル $v(x)$ の作用を受ける。確かめておくべきことは、系において平衡からの非常に大きなずれを考える必要がないという点だけである。これは、すでに用い始めたすべての近似を保つために必要となる。言い換えれば、次の形の Hamiltonian で記述される全体系を研究したいのである。
\begin{equation}
 \begin{aligned}
 H&=H_{\mathrm A}+H_{\mathrm I}+H_{\mathrm B}\\
 &=\frac{p^2}{2M}+v(x)+x\sum_k C_kR_k
 +\sum_k\frac{p_k^2}{2m}+\sum_k\frac12m\omega_k^2R_k^2.
 \end{aligned}
 \tag{3.1}\label{eq:3.1}
\end{equation}
ここで $C_k$ は熱浴の粒子と系 A との結合定数であり、$\omega_k$ は熱浴の粒子の運動の振動数である。式 \eqref{eq:2.16} に現れる作用は、今や
\begin{equation}
 S_{\mathrm B}[\bm R]+S_{\mathrm I}[x,\bm R]
 =\int_0^t L_{\mathrm B}(\bm R,\dot{\bm R})\,dt
 +\int_0^t L_{\mathrm I}(x,\bm R)\,dt
 \tag{3.2}\label{eq:3.2}
\end{equation}
と書ける。ここで
\begin{equation}
 L_{\mathrm B}=\sum_k\frac12m\dot R_k^{,2}
 -\sum_k\frac12m\omega_k^2R_k^2
 \tag{3.3}\label{eq:3.3}
\end{equation}
であり、また
\begin{equation}
 L_{\mathrm I}=-x\sum_k C_kR_k
 \tag{3.4}\label{eq:3.4}
\end{equation}
である。この問題は文献 \cite{ref22,ref23} で厳密に解かれており、式 \eqref{eq:2.16} の影響汎関数の最終的な表式は
\begin{equation}
 \begin{aligned}
 \mathcal F[x,y]=\exp\biggl\{ -\frac1\hbar\int_0^t d\tau\int_0^\tau ds\,
 &[x(\tau)-y(\tau)]\\
 &\times[\alpha(\tau-s)x(s)-\alpha^*(\tau-s)y(s)]\biggr\}
 \end{aligned}
 \tag{3.5}\label{eq:3.5}
\end{equation}
となる。ここで
\begin{equation}
 \alpha(\tau-s)=\sum_k\frac{C_k^2}{2m\omega_k}
 \left[e^{-i\omega_k(\tau-s)}
 +\frac{e^{i\omega_k(\tau-s)}}{e^{\hbar\omega_k/kT}-1}
 +\frac{e^{-i\omega_k(\tau-s)}}{e^{\hbar\omega_k/kT}-1}\right]
 \tag{3.6}\label{eq:3.6}
\end{equation}
である。

% Original PDF page 9 (printed p. 595).
ここでは、環境が初め温度 $T$ の平衡状態にあると仮定した。これは、その初期密度行列を
\begin{equation}
 \rho_{\mathrm B}(\bm R',\bm Q',0)
 =\prod_k\rho_{\mathrm B}^{(k)}(R_k',Q_k',0)
 \tag{3.7}\label{eq:3.7}
\end{equation}
と書けることを意味する。ただし、通常のように
\begin{equation}
 \begin{aligned}
 \rho_{\mathrm B}^{(k)}(R_k',Q_k',0)
 ={}&\frac{m\omega_k}{2\pi\hbar\sinh(\hbar\omega_k/kT)}\\
 &\times\exp\left\{-\frac{m\omega_k}{2\hbar\sinh(\hbar\omega_k/kT)}
 \left[(R_k'^2+Q_k'^2)\cosh\frac{\hbar\omega_k}{kT}-2R_k'Q_k'\right]\right\}.
 \end{aligned}
 \tag{3.8}\label{eq:3.8}
\end{equation}
影響汎関数を表すには、$\alpha(\tau-s)$ の実部と虚部 $\alpha_{\mathrm R}$、$\alpha_{\mathrm I}$ を用いると、はるかに便利である。実際、密度演算子の伝播関数は
\begin{equation}
 \begin{aligned}
 J(x,y,t;x',y',0)={}&\iint\mathcal D x\,\mathcal D y\,
 \exp\biggl\{\frac{i}{\hbar}\biggl(S_{\mathrm A}[x]-S_{\mathrm A}[y]\\
 &\quad-\int_0^t d\tau\int_0^\tau ds\,
 [x(\tau)-y(\tau)]\alpha_{\mathrm I}(\tau-s)[x(s)+y(s)]\biggr)\biggr\}\\
 &\times\exp\biggl\{-\frac1\hbar\int_0^t d\tau\int_0^\tau ds\,
 [x(\tau)-y(\tau)]\alpha_{\mathrm R}(\tau-s)[x(s)-y(s)]\biggr\}
 \end{aligned}
 \tag{3.9}\label{eq:3.9}
\end{equation}
と書ける。ここで $\alpha_{\mathrm R}$ と $\alpha_{\mathrm I}$ は
\begin{equation}
 \alpha_{\mathrm R}(\tau-s)=\sum_k\frac{C_k^2}{2m\omega_k}
 \coth\frac{\hbar\omega_k}{2kT}\cos\omega_k(\tau-s)
 \tag{3.10}\label{eq:3.10}
\end{equation}
および
\begin{equation}
 \alpha_{\mathrm I}(\tau-s)=-\sum_k\frac{C_k^2}{2m\omega_k}
 \sin\omega_k(\tau-s)
 \tag{3.11}\label{eq:3.11}
\end{equation}
で与えられる。

さて、$t=0$ における系の縮約密度行列が与えられれば、式 \eqref{eq:2.14} と \eqref{eq:3.9} によってその時間発展が得られる。もはや熱浴の座標には依存しない。模型 Hamiltonian \eqref{eq:3.1} に関する限り、問題は厳密に解けている。これらすべてを展開した主要な目的は、古典極限で Brownian 粒子を記述することであった。しかし、この段階では、そのようなことが起こるという証拠は何もない。それどころか、この厳密解から不可逆過程を得る方法は見当たらない。

% Original PDF page 10 (printed p. 596).
この問題を解くには、振動子の集団が実際の熱浴として振る舞うためには、幾つかの明確な性質を備えていなければならないことを思い出す必要がある。例えば、対象とする系が失ったエネルギーが有限時間内に系へ戻らないよう、熱浴は無限の大きさを持たなければならない。ここで、振動子の集団が、空間の有限な領域における何らかの物理量の基準振動モードの集合を表すとしよう。例えば、箱の中の電磁場の基準振動モードの集合や、結晶中の原子の平衡位置からの振動であり得る。ただし、振動子の集団についてのより一般的な考察は文献 \cite{ref21} を参照されたい。したがって、体積を無限大にする極限をとるということは、許容される振動子の振動数を連続体として考える必要があることを意味する。

次に、連続体極限で振動子の分布をうまく選べば Brownian 運動に導かれることを示すため、幾つかの議論を行う。その後、選んだ模型を用いて、古典極限で Fokker--Planck 方程式を導く。

式 \eqref{eq:3.1} のように系 A を熱浴 B に結合する代わりに、古典的な外力 $F(\tau)$ を加えて励起すると仮定しよう。この場合、密度行列の伝播関数が
\begin{equation}
 \begin{aligned}
 J(x,y,t;x',y',0)=\iint\mathcal D x\,\mathcal D y\,
 \exp\left\{\frac{i}{\hbar}\left(S_{\mathrm A}[x]-S_{\mathrm A}[y]
 +\int_0^t[x(\tau)-y(\tau)]F(\tau)\,d\tau\right)\right\}
 \end{aligned}
 \tag{3.12}\label{eq:3.12}
\end{equation}
で与えられることは容易に理解できる\cite{ref22}。

さらに、$F(\tau)$ 自体は正確には分からず、異なる履歴 $F(\tau)$ の確率分布汎関数 $P[F(\tau)]$ が分かっているものとすると、平均された密度行列の伝播関数は
\begin{equation}
 \begin{aligned}
 J(x,y,t;x',y',0)={}&\iiint\mathcal D x\,\mathcal D y\,\mathcal D F\,P[F(\tau)]\\
 &\times\exp\left\{\frac{i}{\hbar}\left(S_{\mathrm A}[x]-S_{\mathrm A}[y]
 +\int_0^t[x(\tau)-y(\tau)]F(\tau)\,d\tau\right)\right\}
 \end{aligned}
 \tag{3.13}\label{eq:3.13}
\end{equation}
となる。式 \eqref{eq:3.13} における $F(\tau)$ についての汎関数積分は、よく知られた特性汎関数\cite{ref22}
\begin{equation}
 \begin{aligned}
 \phi\left[\frac{x(\tau)-y(\tau)}{\hbar}\right]
 \equiv\int\mathcal D F(\tau)\,P[F(\tau)]
 \exp\left\{i\int_0^t\left[\frac{x(\tau)-y(\tau)}{\hbar}\right]F(\tau)\,d\tau\right\}
 \end{aligned}
 \tag{3.14}\label{eq:3.14}
\end{equation}
である。式 \eqref{eq:3.14} を用いると、異なる二つの時刻における力の相関を次のように求めることができる。
% The words "instants as" at the start of original PDF page 11 are included in the preceding sentence.

% Source PDF p.11 (printed p.597); lead-in to (3.15) is in part1.tex.
\begin{equation}
\langle F(\xi)F(s)\rangle
=-\left.
\frac{\delta^2\phi}
{\delta\!\left[\dfrac{x(\xi)-y(\xi)}{\hbar}\right]
 \delta\!\left[\dfrac{x(s)-y(s)}{\hbar}\right]}
\right|_{x=y},
\tag{3.15}\label{eq:3.15}
\end{equation}
ここで $\delta$ は汎関数微分（または変分）を表す。$P[F(\tau)]$ が Gaussian 分布であるとき、$\phi$ は非常に簡単な形をとる\cite{ref22}。
\begin{equation}
\begin{aligned}
\phi\!\left[\frac{x(\tau)-y(\tau)}{\hbar}\right]
=\exp\biggl\{-\frac{1}{\hbar^2}\int_0^t\!\mathrm d\tau\int_0^\tau\!\mathrm ds\,
[x(\tau)-y(\tau)]A(\tau-s)[x(s)-y(s)]\biggr\},
\end{aligned}
\tag{3.16}\label{eq:3.16}
\end{equation}
これと\eqref{eq:3.15}から、
\begin{equation}
\langle F(\xi)F(s)\rangle=A(\xi-s)
\tag{3.17}\label{eq:3.17}
\end{equation}
を得る。

ここで\eqref{eq:3.14}を\eqref{eq:3.13}に代入し、$P[F(\tau)]$ が Gaussian 型の汎関数分布であると仮定すると、伝播関数 $J$ は
\begin{equation}
\begin{aligned}
J(x,y,t;x',y',0)
={}&\iint\mathcal D x\,\mathcal D y\,
\exp\!\left\{\frac{\mathrm i}{\hbar}(S_A[x]-S_A[y])\right\}\\
&\times\exp\biggl\{-\frac{1}{\hbar^2}
\int_0^t\!\mathrm d\tau\int_0^\tau\!\mathrm ds\,
[x(\tau)-y(\tau)]A(\tau-s)[x(s)-y(s)]\biggr\}
\end{aligned}
\tag{3.18}\label{eq:3.18}
\end{equation}
となる。

\eqref{eq:3.18}と\eqref{eq:3.9}を比較すると、\eqref{eq:3.9}の指数に現れる付加的な虚部を除いて、両者は非常によく似ている。この違いにかかわらず、指数の実部は二つの場合で同じ役割を果たすと考えられる。\eqref{eq:3.18}の $A(\tau-s)$ と\eqref{eq:3.9}の $\hbar\alpha_R(\tau-s)$ の物理的意味が異なるはずはない。したがって、$\hbar\alpha_R(\tau-s)$ は古典領域における力の相関を与えなければならない。ここでの問題は、どのような条件の下で、この量子力学的模型が古典的な Brownian 粒子に作用する確率的な力の相関を再現するか、ということである。これに答えるのは難しくない。高温での $\hbar\alpha_R(\tau-s)$ の振る舞いを調べればよい。$kT\gg\hbar\omega_k$ のとき、\eqref{eq:3.10}は
\begin{equation}
\alpha_R(\tau-s)\approx
\frac{kT}{m\hbar}\sum_i\frac{C_i^2}{\omega_i^2}\cos\omega_i(\tau-s)
+\frac{\hbar}{12mkT}\sum_i C_i^2\cos\omega_i(\tau-s)+\cdots
\tag{3.19}\label{eq:3.19}
\end{equation}
となる。求めているのは $\hbar\alpha_R(\tau-s)$ の展開なので、その主要項は
\begin{equation}
\hbar\alpha_R(\tau-s)\approx
\frac{kT}{m}\sum_i\frac{C_i^2}{\omega_i^2}\cos\omega_i(\tau-s)
+\mathcal O(\hbar^2)
\tag{3.20}\label{eq:3.20}
\end{equation}
である。これが、この量子力学的模型から得られる古典的な相関である。

% Source PDF p.12 (printed p.598).
\eqref{eq:3.20}を、Brownian 運動の古典論が与えるよく知られた力の相関\cite{ref01}
\begin{equation}
\langle F(\tau)F(s)\rangle=2\eta kT\delta(\tau-s)
\tag{3.21}\label{eq:3.21}
\end{equation}
と比較したい。ここで $\eta$ は減衰定数である。そのため、密度 $\rho_D(\omega)$ の連続な振動子分布を考える。すると、
\begin{equation}
\hbar\alpha_R(\tau-s)\approx\frac{kT}{m}
\int_0^\infty\rho_D(\omega)\frac{C^2(\omega)}{\omega^2}
\cos\omega(\tau-s)\,\mathrm d\omega
\tag{3.22}\label{eq:3.22}
\end{equation}
となる。ここで、次のように選ぶ\cite{ref20}。
\begin{equation}
\rho_D(\omega)C^2(\omega)=
\begin{cases}
\dfrac{2m\eta\omega^2}{\pi},&\omega<\Omega,\\[0.5em]
0,&\omega>\Omega.
\end{cases}
\tag{3.23}\label{eq:3.23}
\end{equation}
振動子分布に高周波数カットオフ $\Omega$ を導入すると、\eqref{eq:3.22}は
\begin{equation}
\hbar\alpha_R(\tau-s)=\langle F(\tau)F(s)\rangle
=2\eta kT\frac{1}{\pi}\frac{\sin\Omega(\tau-s)}{\tau-s}
\tag{3.24}\label{eq:3.24}
\end{equation}
となり、$\Omega\to\infty$ とすれば\eqref{eq:3.21}に近づく。言い換えれば、特徴的な時間 $\Omega^{-1}$ よりもはるかに長い時間に着目するとき、\eqref{eq:3.21}が再現される。これは、この場合には\eqref{eq:3.23}の低周波数での振る舞いが重要であることを意味する。この事実は Brownian 運動の古典論と整合している。なぜなら、\eqref{eq:3.21}が成り立つのは、熱浴の特徴的な緩和時間よりも長い時間を考える場合に限られるからである。

\eqref{eq:3.10}を
\begin{equation}
\hbar\alpha_R(\tau-s)=\frac{\hbar}{\pi}\int_0^\infty
\mathrm d\omega\,\chi''_{FF}(\omega)
\coth\frac{\hbar\omega}{2kT}\cos\omega(\tau-s)
\tag{3.25}\label{eq:3.25}
\end{equation}
と書けば、さらに少し議論を進められる。ここでは
\begin{equation}
\chi''_{FF}(\omega)=\sum_i\frac{\pi C_i^2}{2m\omega_i}
\delta(\omega-\omega_i)
\tag{3.26}\label{eq:3.26}
\end{equation}
と定義した。これにより、$\hbar\alpha_R(\tau-s)$ を、異なる時刻におけるある演算子 $\widetilde F(t)$ の時間相関関数とみなすことができる。したがって、
\begin{equation}
\hbar\alpha_R(\tau-s)=
\tfrac12\langle\{\widetilde F(\tau),\widetilde F(s)\}\rangle
-\langle\widetilde F(\tau)\rangle\langle\widetilde F(s)\rangle
\tag{3.27}\label{eq:3.27}
\end{equation}
と定義すれば、古典極限で $\widetilde F(t)$ が古典的な揺動力 $F(t)$ になる限り、この極限で\eqref{eq:3.21}が再現される。式\eqref{eq:3.25}はよく知られた揺動散逸定理である\cite{ref24}。

% Source PDF p.13 (printed p.599).
\eqref{eq:3.25}と\eqref{eq:3.26}を用いる利点は、\eqref{eq:3.21}を再現する適切な振動子分布を見いだせるだけでなく、$\chi''_{FF}(\omega)$ の和則を用いて、その分布に対する補正も求められることである。もちろん、これらの補正が重要になるのは、系の短時間の振る舞いに着目するときに限られる\cite{ref24}。

縮約密度演算子の伝播関数の最終的な形を得るには、先に述べた指数中の付加的な虚部を求める必要がある。この段階では、振動子分布はすでに\eqref{eq:3.23}によって定まっているので、それは非常に簡単である。この式を\eqref{eq:3.11}に代入すると、
\begin{equation}
\alpha_I(\tau-s)=\frac{\eta}{2\pi}\frac{\mathrm d}{\mathrm d(\tau-s)}
\int_{-\Omega}^{\Omega}\cos\omega(\tau-s)\,\mathrm d\omega
\tag{3.28}\label{eq:3.28}
\end{equation}
を得る。これは $\Omega\to\infty$ のとき
\begin{equation}
\alpha_I(\tau-s)=\eta\frac{\mathrm d}{\mathrm d(\tau-s)}\delta(\tau-s)
\tag{3.29}\label{eq:3.29}
\end{equation}
に近づく。最後に、\eqref{eq:3.25}と\eqref{eq:3.28}を\eqref{eq:3.9}に用いると、
\begin{equation}
\begin{aligned}
J(x,y,t;x',y',0)
={}&\iint\mathcal D x\,\mathcal D y\,
\exp\biggl\{\frac{\mathrm i}{\hbar}\biggl(
S_A[x]-S_A[y]\\
&-\int_0^t\!\mathrm d\tau\int_0^\tau\!\mathrm ds\,
\eta[x(\tau)-y(\tau)]
\frac{\mathrm d}{\mathrm d(\tau-s)}\delta(\tau-s)[x(s)+y(s)]
\biggr)\biggr\}\\
&\times\exp\biggl\{-\frac{1}{\hbar}\int_0^\Omega
\frac{\eta\omega}{\pi}\coth\frac{\hbar\omega}{2kT}\,
\mathrm d\omega\int_0^t\!\mathrm d\tau\int_0^\tau\!\mathrm ds\,
[x(\tau)-y(\tau)]\\
&\hspace{9em}\times\cos\omega(\tau-s)[x(s)-y(s)]\biggr\}
\end{aligned}
\tag{3.30}\label{eq:3.30}
\end{equation}
となる。この式は、デルタ関数の微分を含む項を部分積分することで、さらに簡単にできる。計算は直接的であり、結果は
\begin{equation}
\begin{aligned}
&\int_0^t\!\mathrm d\tau\int_0^\tau\!\mathrm ds\,
\eta[x(\tau)-y(\tau)]\frac{\mathrm d}{\mathrm d(\tau-s)}
\delta(\tau-s)[x(s)+y(s)]\\
&\quad=-\int_0^t\eta[x^2(\tau)-y^2(\tau)]\delta(0)\,\mathrm d\tau\\
&\qquad\quad+\frac{\eta}{2}\int_0^t
[x(\tau)\dot x(\tau)+x(\tau)\dot y(\tau)
-y(\tau)\dot x(\tau)-y(\tau)\dot y(\tau)]\,\mathrm d\tau
\end{aligned}
\tag{3.31}\label{eq:3.31}
\end{equation}
となる。

% Source PDF p.14 (printed p.600).
一見すると、$\delta(0)$ による発散が気になるだろう。しかし、この項の解釈には注意が必要である。これは $\Omega\to\infty$ のときの式\eqref{eq:3.28}から生じたものであることを思い出すとよい。したがって、
\begin{equation}
\delta(0)=\lim_{\Omega\to\infty}
\left.\frac{\sin\Omega(\tau-s)}{\pi(\tau-s)}\right|_{\tau=s}
=\lim_{\Omega\to\infty}\frac{\Omega}{\pi}
\tag{3.32}\label{eq:3.32}
\end{equation}
と書ける。\eqref{eq:3.31}から分かるように、$\delta(0)$ は減衰定数 $\eta$ と掛け合わされて現れるので、積 $\eta\delta(0)$ は実際には
\begin{equation}
\eta\delta(0)=\lim_{\Omega\to\infty}\frac{\eta\Omega}{\pi}
\tag{3.33}\label{eq:3.33}
\end{equation}
である。

これは非常に微妙な点である。一見すると、$\Omega\to\infty$ のときこの式は当然発散すると言うべきだろう。一方、この主張に対して、極限 $\Omega\to\infty$ の意味は、Brownian 粒子の運動の特徴的な周波数よりも $\Omega$ をはるかに大きくとることなのだ、と反論することもできる。したがって、$\eta$ と $\Omega$ の関係によっては、積 $\eta\Omega$ が有限になると期待できる（$\eta$ が本当に現象論的な減衰定数であるならば）。しかし、この議論にも批判の余地はある。

何が起こるかをより形式的に調べるには、\eqref{eq:3.23}の両辺を積分すればよい。熱浴中に $N$ 個の振動子があり、$C(\omega)=C$ が定数であると仮定すると、\eqref{eq:3.33}の右辺は
\begin{equation}
\frac{\eta\Omega}{\pi}=\frac{3NC^2}{2m\Omega^2}
\tag{3.34}\label{eq:3.34}
\end{equation}
と表せる。この式では $N$ と $\Omega$ がともに無限大に近づく。したがって、$\eta\Omega$ が発散するかどうかは $N\Omega^{-2}$ に依存する。$\Omega\propto N^{1/2}$ となる振動子集団を模型として採用すれば、\eqref{eq:3.34}は明らかに収束する。

しかし、$\alpha_I(\tau-s)$ を含む積分を評価した後で初めて大きな $\Omega$ の極限を調べれば、これらの議論はすべて避けられる。そうすると、\eqref{eq:3.31}の右辺の第1積分は
\begin{equation}
\begin{aligned}
-\frac{\eta\Omega}{\pi}\int_0^t[x^2(\tau)-y^2(\tau)]\,\mathrm d\tau
+\frac{\eta[x(0)+y(0)]}{\pi}
\int_0^t\frac{\sin\Omega\tau}{\tau}[x(\tau)-y(\tau)]\,\mathrm d\tau
\end{aligned}
\tag{3.35}\label{eq:3.35}
\end{equation}
となる。ここでは $\delta(\tau-s)$ を適切に $(1/\pi)\sin[\Omega(\tau-s)]/(\tau-s)$ で置き換えた。大きな $\Omega$ の極限は、時間 $t\approx\omega^{-1}\gg\Omega^{-1}$ に着目しているという事実を反映しているにすぎない。その結果、\eqref{eq:3.35}の第2項は $\frac12\eta[x^2(0)-y^2(0)]$ となり、これは明らかに第1項よりもはるかに小さい。この結果により、極限 $\Omega\to\infty$ をとることなく、\eqref{eq:3.31}において $\eta\delta(0)$ を $\eta\Omega/\pi$ で置き換えることができる。周波数シフト（後述の定義\eqref{eq:3.37}を参照）に関する詳しい考察については、本論文の第4節および文献\cite{ref21}を参照されたい。

% Source PDF p.15 (printed p.601).
緩和定数 $\gamma$ を
\begin{equation}
\gamma\equiv\eta/2M
\tag{3.36}\label{eq:3.36}
\end{equation}
と定義し、周波数シフト $\Delta\omega$ を
\begin{equation}
(\Delta\omega)^2\equiv\frac{4\gamma\Omega}{\pi}
\tag{3.37}\label{eq:3.37}
\end{equation}
と定義すると、最終的に\eqref{eq:3.30}は
\begin{equation}
\begin{aligned}
J(x,y,t;x',y',0)
={}&\iint\mathcal D x\,\mathcal D y\,
\exp\biggl\{\frac{\mathrm i}{\hbar}\biggl(
S_R[x]-S_R[y]\\
&\hspace{4em}-M\gamma\int_0^t
(x\dot x-y\dot y+x\dot y-y\dot x)\,\mathrm d\tau\biggr)\biggr\}\\
&\times\exp\biggl\{-\frac{1}{\hbar}\frac{2M\gamma}{\pi}
\int_0^\Omega\!\mathrm d\omega\,\omega\coth\frac{\hbar\omega}{2kT}
\int_0^t\!\mathrm d\tau\int_0^\tau\!\mathrm ds\,
[x(\tau)-y(\tau)]\\
&\hspace{9em}\times\cos\omega(\tau-s)[x(s)-y(s)]\biggr\}
\end{aligned}
\tag{3.38}\label{eq:3.38}
\end{equation}
と書ける。ここでは、繰り込まれた作用 $S_R$ を
\begin{equation}
S_R=\int_0^t\bigl[\tfrac12 M\dot x^2-v(x)\bigr]\,\mathrm d\tau
+\int_0^t\tfrac12 M(\Delta\omega)^2x^2\,\mathrm d\tau
\tag{3.39}\label{eq:3.39}
\end{equation}
で導入した。

言い換えると、$S_R$ は、周波数 $\Delta\omega$ の調和項を差し引くことによってポテンシャル $v(x)$ を繰り込んだ作用である。以後、これを繰り込まれたポテンシャル $v_R(x)$ と呼ぶ。$v(x)=\frac12 M\omega^2x^2$（単純調和振動子）の場合には、繰り込まれた周波数 $\omega_R$ も $\omega_R^2\equiv\omega^2-(\Delta\omega)^2$ で定義する。粒子が感じるこのような有効ポテンシャルは、粒子を熱浴に結合すると一般に生じ得る\cite{ref17}。ここでは、環境に用いた模型のために、非常に簡単な形をとっている。

ここで、われわれの式\eqref{eq:3.38}と、これまでに導かれた他の式との関連をいくつか強調しておきたい。まず、この式の指数の虚部が、Morse と Feshbach\cite{ref25}の提案した二変数の作用に非常によく似ていることに注意しよう。簡単のため $v_R(x)=0$ とし、この指数の虚部を $S_{\mathrm{eff}}$ と呼ぶと、
\begin{equation}
S_{\mathrm{eff}}=\int_0^t\bigl[\tfrac12 M\dot x^2-\tfrac12 M\dot y^2
-M\gamma(x\dot x-y\dot y+x\dot y-y\dot x)\bigr]\,\mathrm d\tau
\tag{3.40}\label{eq:3.40}
\end{equation}
となる。元の変数を $\pi/4$ 回転して得られる新しい変数の組 $\{\bar x,\bar y\}$ を定義すると、

% Source PDF p.16 (printed p.602).
\begin{equation}
S_{\mathrm{eff}}=\int_0^t\left[M\dot{\bar x}\dot{\bar y}
-\frac{\eta}{2}(\bar x\dot{\bar y}-\bar y\dot{\bar x})\right]\mathrm d\tau
+\eta\bigl(\bar x(t)\bar y(t)-\bar x(0)\bar y(0)\bigr)
\tag{3.41}\label{eq:3.41}
\end{equation}
を得る。ここでは、Morse と Feshbach の Lagrangian が、時間について積分された形で自然に現れる。したがって、われわれの式\eqref{eq:3.38}は、その Lagrangian を正しく解釈する手掛かりを与えている。変数 $\bar x$ と $\bar y$（または $x$ と $y$）を、互いに相互作用する二つの異なる系を表すものと考えることはできない。これらは、時間を順方向および逆方向に進む、結合した経路を表す。この二つの経路は、$\widetilde\rho(x,y,t)$ の時間発展を記述するために必要である。熱浴との一般的な相互作用に対するこれらの経路の結合の図式展開は、Kjeldysh 形式\cite{ref26}を用いて文献\cite{ref17}で行われている。

数年前、Nemes と Piza\cite{ref27}は、Morse と Feshbach の Lagrangian に基づいて、量子力学における減衰を扱う現象論的理論を構築した。そこでは $\bar y$ を非物理的な変数とみなし、それについてトレースをとっている。この手続きは非常に興味深いが、最終的には、長時間で波束の幅が凍結するという、線形応答理論に反する結果に至っている。

われわれの議論により近い文脈で、Papadopoulos\cite{ref28}は\eqref{eq:3.38}に非常によく似た式を導いた。しかし、両者の式には二つの大きな違いがある。第一の違いは\eqref{eq:3.38}の実指数に関するもので、彼は周波数 $\omega$ の調和振動子の場合に、それを
\begin{equation}
\text{（実指数）}=-\frac{1}{\hbar^2}M\gamma\hbar\omega
\coth\frac{\hbar\omega}{2kT}\int_0^t[x(\tau)-y(\tau)]^2\,\mathrm d\tau
\tag{3.42}\label{eq:3.42}
\end{equation}
と書いている。これは、われわれの式の特殊な場合であることが示せる。この点は後に第6節で詳しく論じる。第二の違いは、彼の式には\eqref{eq:3.38}の虚指数にある $\gamma$ に依存する項がないことである。その種の項の最終的な効果を説明するため、彼は時間に陽に依存する減衰をその場しのぎで導入している。しかし、$\gamma$ に依存する項は微視的な計算から自然に生じるので、そのような導入はまったく不要であるとわれわれは考える。

\sectionlink{速度への結合}
\label{sec:4}
全 Lagrangian が
\begin{equation}
L(x,\dot x,\boldsymbol q,\dot{\boldsymbol q})
=\tfrac12 M\dot x^2-v(x)+x\sum_i e_i\dot q_i
+\sum_i\tfrac12 m\dot q_i^2-\sum_i\tfrac12 m\omega_i^2q_i^2
\tag{4.1}\label{eq:4.1}
\end{equation}
で与えられるように、熱浴に結合した粒子の運動を調べたいとしよう。

% Source PDF p.17 (printed p.603).
この式がこれまで扱ってきたものと異なるのは、結合項だけである（\eqref{eq:3.4}と\eqref{eq:4.1}を比較せよ）。

したがって、少なくとも直接には、前節の結果をこの新しい模型に適用することはできない。これから示すのは、\eqref{eq:4.1}を、座標に結合する非常によく似た Lagrangian へ変換できるということである。

この問題を克服するため、次の新しい座標の組を導入しよう。
\begin{equation}
y_i=m\dot q_i+e_i x
\tag{4.2}\label{eq:4.2}
\end{equation}
$L(x,\dot x,\boldsymbol y,\dot{\boldsymbol y})$ と書いて Euler--Lagrange 方程式を用いても、\eqref{eq:4.1}から得られる正しい運動方程式は再現されない。これは、$L$ が変換\eqref{eq:4.2}の下で不変ではないことを意味する。新しい変数の組における正しい Lagrangian が
\begin{equation}
\begin{aligned}
\widetilde L(x,\dot x,\boldsymbol y,\dot{\boldsymbol y})
={}&\tfrac12 M\dot x^2-v(x)-\sum_i\frac{e_i^2x^2}{2m}
+\sum_i\frac{e_i y_i x}{m}\\
&+\tfrac12\sum_i\frac{\dot y_i^2}{m\omega_i^2}
-\tfrac12\sum_i\frac{y_i^2}{m}
\end{aligned}
\tag{4.3}\label{eq:4.3}
\end{equation}
となることは、非常に簡単な計算で確かめられる。さらに変数 $R_i$ を
\begin{equation}
R_i\equiv\frac{y_i}{m\omega_i}
\tag{4.4}\label{eq:4.4}
\end{equation}
で定義すれば、\eqref{eq:4.3}を
\begin{equation}
\begin{aligned}
\widetilde L(x,\dot x,\boldsymbol R,\dot{\boldsymbol R})
={}&\tfrac12 M\dot x^2-v(x)-\sum_i\frac{C_i^2x^2}{2m\omega_i^2}
+\sum_i C_i xR_i\\
&+\sum_i\tfrac12 m\dot R_i^2-\sum_i\tfrac12 m\omega_i^2R_i^2
\end{aligned}
\tag{4.5}\label{eq:4.5}
\end{equation}
と書き直せる。ここで $C_i\equiv e_i\omega_i$ であり、これは前節で用いた Lagrangian とまったく同じ形である。ここでの違いは、変換\eqref{eq:4.2}によって裸のポテンシャル $v(x)$ に新たな二次の項が導入されたことである。したがって、新しいポテンシャルは
\begin{equation}
\bar v(x)=v(x)+\sum_i\frac{C_i^2}{2m\omega_i^2}x^2
\tag{4.6}\label{eq:4.6}
\end{equation}
となる。

これで、以前と同じように議論を進められる。問題は第3節のものに帰着された。縮約密度演算子の伝播関数は、再び式\eqref{eq:3.38}で与えられる。ここで注目すべきなのは、繰り込まれたポテンシャル $\bar v_R(x)$ の形であり、これは
\begin{equation}
\bar v_R(x)=\bar v(x)-\tfrac12 M(\Delta\omega)^2x^2
\tag{4.7}\label{eq:4.7}
\end{equation}
と書ける。

周波数シフト $\Delta\omega$ は式\eqref{eq:3.37}で定義されているが、それを得る際には、すでに特定の振動子分布を仮定していた。$\alpha_I(\tau-s)$ を、より一般的な式\eqref{eq:3.11}で与えられるものとして、\eqref{eq:3.28}から\eqref{eq:3.31}までの手順をたどれば、より一般的に議論できる。その場合、$\Delta\omega$ は

% Source PDF p.18 (printed p.604).
\begin{equation}
M(\Delta\omega)^2
=\left.2\sum_i\frac{C_i^2}{2m\omega_i^2}
\cos\omega_i(\tau-s)\right|_{\tau=s}
=2\sum_i\frac{C_i^2}{2m\omega_i^2}
\tag{4.8}\label{eq:4.8}
\end{equation}
を満たす。分布\eqref{eq:3.23}が成り立つと仮定すれば、これは明らかに\eqref{eq:3.37}に帰着する。\eqref{eq:4.8}を用いる利点は、$\Delta\omega$ に依存する補正ポテンシャルが、\eqref{eq:4.6}で加わった部分を正確に打ち消すことを示せる点にある。さらに、この事実は、用いている振動子分布にまったく依存しない。したがって、\eqref{eq:4.1}で提案した模型については、繰り込まれたポテンシャルは裸のポテンシャルと同じであると言える。すなわち、
\begin{equation}
\bar v_R(x)=v(x)
\tag{4.9}\label{eq:4.9}
\end{equation}
である。

この節で論じた一見するとかなり抽象的な点の物理的な意味は、環境に散逸的に結合しているすべての系が、その結合による周波数シフト（またはポテンシャルの繰り込み）を必ずしも受けるわけではない、という結論にある。この点は、LCR 回路および類似の系の量子化においてある程度重要であり（実際、この点を理解しなかったために、最近の文献で不要な混乱が生じた）、文献\cite{ref21}でさらに論じられている。

\sectionlink{Fokker--Planck 方程式}
\label{sec:5}
前の二つの節で、ある特定の性質を持つ熱浴と相互作用する粒子の縮約密度演算子について、伝播関数の式を導いた。また、古典極限では、\eqref{eq:3.38}の実指数因子が、Brownian 粒子に作用する確率的な力の特性汎関数に帰着することも示した。同時に、被積分関数の付加的な虚部、すなわち\eqref{eq:3.38}の $x\dot x$ などを含む項にも、明確な形を与えた。ここで答えたい問題は、この項が意味を持つかどうかである。言い換えれば、この項は粒子の古典的な Brownian 運動と両立するのだろうか。この問いに答えるため、まず半古典領域における縮約密度演算子の運動方程式を調べる。この方程式を古典的な Brownian 粒子のものと比較するには、位相空間表示で書く方法を見つけなければならない。古典物理における位相空間分布の運動方程式と比較するには、それしか方法がない。Hilbert 空間から位相空間への変換は、いわゆる Wigner 分布関数\cite{ref29}を用いて行える。Wigner の形式に関する最近の総説については、文献\cite{ref03,ref30}を参照されたい。

$2kT\gtrsim\hbar\Omega\gg\hbar\omega_R$ の場合に\eqref{eq:3.38}を書き下すことから始めよう。ここで $\Omega$ は熱浴の振動子のカットオフ周波数である。すると、

% Source PDF p.19 (printed p.605).
\begin{equation}
\begin{aligned}
J(x,y,t;x',y',0)
={}&\iint\mathcal D x\,\mathcal D y\,
\exp\biggl\{\frac{\mathrm i}{\hbar}\biggl(
S_R[x]-S_R[y]\\
&\hspace{4em}-M\gamma\int_0^t
[x\dot x-y\dot y+x\dot y-y\dot x]\,\mathrm d\tau\biggr)\biggr\}\\
&\times\exp\biggl\{-\frac{2M\gamma kT}{\hbar^2}
\int_0^t[x(\tau)-y(\tau)]^2\,\mathrm d\tau\biggr\}
\end{aligned}
\tag{5.1}\label{eq:5.1}
\end{equation}
となる。

ここで、この高温における量子力学的な伝播関数の意味を気にするべきであろう。もちろん、古典極限 $kT\gg\hbar\omega_R$ を考えながら、その量子的な形をなお保つならば、この式は無意味になる。しかし、これは古典極限における縮約密度演算子の運動方程式を得るための最初の段階であることを、忘れてはならない。後に、この近似と整合するよう適切な措置をとる。\eqref{eq:3.38}ではなく\eqref{eq:5.1}を用いるのは、これから現れる式を簡単にするためだけである。

これから行うのは、Feynman が文献\cite{ref22}で汎関数積分形式から Schrödinger 方程式を導いたときの手続きに従うことである。ここでの問題は、その問題と類似している。文献\cite{ref22}では、波動関数の伝播関数から $\psi$ の運動方程式を得るのに対し、ここでは、密度演算子の伝播関数から $\widetilde\rho$ のマスター方程式を得ようとしている。

時刻 $t$ での縮約密度演算子が分かっていて、$\epsilon\to0$ として、$t+\epsilon$ での値を求めたいとしよう。\eqref{eq:2.14}から、
\begin{equation}
\widetilde\rho(x,y,t+\epsilon)=\iint\mathrm dx'\,\mathrm dy'\,
J(x,y,t+\epsilon;x',y',t)\widetilde\rho(x',y',t)
\tag{5.2}\label{eq:5.2}
\end{equation}
となる。

\eqref{eq:5.2}の伝播関数 $J$ は、$\epsilon$ が小さいとき、非常に簡単な形に書ける。そのためには、短い時間間隔において、任意の正則な経路は直線で近似できることを思い出すだけでよい。したがって、短い時間間隔内の経路に関する汎関数積分は、被積分関数の値に規格化定数を掛けたものに等しいと置ける\cite{ref22}。すると、
\begin{equation}
\begin{aligned}
J(x,y,t+\epsilon;x',y',t)
\approx{}&\frac{1}{A^2}\exp\biggl\{\frac{\mathrm i}{\hbar}\biggl[
\int_t^{t+\epsilon}\bigl(\tfrac12 M\dot x^2-v_R(x)\bigr)\,\mathrm d\tau\\
&-\int_t^{t+\epsilon}\bigl(\tfrac12 M\dot y^2-v_R(y)\bigr)\,\mathrm d\tau\\
&-\int_t^{t+\epsilon}M\gamma
[x\dot x-y\dot y+y\dot x-x\dot y]\,\mathrm d\tau\biggr]\biggr\}\\
&\times\exp\biggl\{-\frac{1}{\hbar^2}\int_t^{t+\epsilon}
2M\gamma kT[x-y]^2\,\mathrm d\tau\biggr\}
\end{aligned}
\tag{5.3}\label{eq:5.3}
\end{equation}
となる。

% Source PDF p.20 (printed p.606).
上に現れたすべての積分も、$\epsilon\to0$ のとき近似できる。$x(t+\epsilon)=x$、$y(t+\epsilon)=y$、$x(t)=x'$、$y(t)=y'$ とし、$J$ の式で
\begin{equation}
\dot x\approx\frac{x-x'}{\epsilon},\qquad
\dot y\approx\frac{y-y'}{\epsilon},\qquad
\int_t^{t+\epsilon}f(x(\tau))\,\mathrm d\tau
\approx\epsilon f\!\left(\frac{x+x'}{2}\right)
\tag{5.4}\label{eq:5.4}
\end{equation}
を用いると、式\eqref{eq:5.2}は
\begin{equation}
\begin{aligned}
\widetilde\rho(x,y,t+\epsilon)
={}&\int_{-\infty}^{\infty}\!\int_{-\infty}^{\infty}
\mathrm d\beta_1\,\mathrm d\beta_2\,
\exp\biggl\{\frac{\mathrm iM\beta_1^2}{2\epsilon\hbar}
-\frac{\mathrm i\epsilon}{\hbar}v_R\!\left(x-\frac{\beta_1}{2}\right)
-\frac{\mathrm iM\beta_2^2}{2\epsilon\hbar}\\
&+\frac{\mathrm i\epsilon}{\hbar}v_R\!\left(y-\frac{\beta_2}{2}\right)
-\frac{\mathrm iM\gamma}{\hbar}\left(x-\frac{\beta_1}{2}\right)\beta_2
+\frac{\mathrm iM\gamma}{\hbar}\left(y-\frac{\beta_2}{2}\right)\beta_1\\
&-\frac{\mathrm iM\gamma}{\hbar}\left(x-\frac{\beta_1}{2}\right)\beta_1
+\frac{\mathrm iM\gamma}{\hbar}\left(y-\frac{\beta_2}{2}\right)\beta_2
-\frac{2M\gamma kT\epsilon}{\hbar^2}(x-y)^2\\
&-\frac{2M\gamma kT\epsilon}{\hbar^2}(x-y)(\beta_1-\beta_2)
-\frac{M\gamma kT\epsilon}{\hbar^2}(\beta_1-\beta_2)^2\biggr\}\\
&\times\widetilde\rho(x-\beta_1,y-\beta_2,t)
\end{aligned}
\tag{5.5}\label{eq:5.5}
\end{equation}
となる。ここで $x-x'\equiv\beta_1$ および $y-y'\equiv\beta_2$ とした。これで、$\epsilon\to0$ の極限で\eqref{eq:5.5}を評価できる。

まず、\eqref{eq:5.5}の積分には、指数が $\epsilon^{-1}$ に比例する、非常に速く振動する二つの項が含まれていることに注意しなければならない。主要な寄与が、$\beta_1$ と $\beta_2$ が非常に小さいところから来ることは明らかである。そうでなければ、因子 $\exp[(\mathrm iM/2\epsilon\hbar)(\beta_1^2-\beta_2^2)]$ は激しく振動し、\eqref{eq:5.5}に有限の寄与を与えない。より具体的には、
\begin{equation}
\beta_1\approx\beta_2\approx\left(\frac{\epsilon\hbar}{M}\right)^{1/2}
\tag{5.6}\label{eq:5.6}
\end{equation}
という領域が必要である。この領域では、両方の指数因子の位相が1の程度だけ変化するからである。ここで、有限の $\beta$ を持つ領域 $\beta_1\approx\beta_2\approx\beta$ からの寄与はどうなるのか、と疑問に思うかもしれない。この場合、二つの指数因子を合わせた位相が1の程度だけ変化するのは、
\begin{equation}
\Delta\beta\equiv\beta_1-\beta_2\approx\frac{\epsilon\hbar}{M\beta}
\tag{5.7}\label{eq:5.7}
\end{equation}
のときである。

ここで、新しい変数 $\beta'_1\equiv\beta_1-\gamma(x-y)\epsilon$ および $\beta'_2\equiv\beta_2+\gamma(x-y)\epsilon$ を定義し、\eqref{eq:5.5}の指数因子を $\epsilon$ について展開すると、$\epsilon\to0$ の極限で、被積分関数中の $\Delta\beta$ に依存するすべての項が $\mathcal O(\epsilon^2)$ になることを容易に示せる。これは、領域 $\beta_1\approx\beta_2\approx\beta$ を無視してよいことを意味する。

次に行う手順は簡単である。$\beta_1\approx\beta_2\approx0$ のまわりで $\widetilde\rho(x-\beta_1,y-\beta_2,t)$ を展開し、積の中で $\mathcal O(\epsilon)$ の項をすべて残そう。$\beta'_1$ および $\beta'_2$ で表すと、次のようになる。
% The final lead-in includes the first sentence continuation on source p.21.

% 原文 PDF p.21（誌面 p.607）。直前の導入文は part2.tex に収録。
\begin{equation}
\begin{aligned}
\tilde\rho(x,y,t)={}&\iint\frac{\mathrm d\beta'_1\,\mathrm d\beta'_2}{A^2}
\exp\!\left(\frac{iM{\beta'_1}^2}{2\epsilon\hbar}\right)
\exp\!\left(-\frac{iM{\beta'_2}^2}{2\epsilon\hbar}\right)
\Bigg[\tilde\rho(x,y,t)
-\frac{\partial\tilde\rho}{\partial x}\beta'_1
-\frac{\partial\tilde\rho}{\partial y}\beta'_2\\
&-\frac{\partial\tilde\rho}{\partial x}\gamma(x-y)\epsilon
+\frac{\partial\tilde\rho}{\partial y}\gamma(x-y)\epsilon
+\frac12\frac{\partial^2\tilde\rho}{\partial x^2}{\beta'_1}^2
+\frac{\partial^2\tilde\rho}{\partial x\partial y}\beta'_1\beta'_2\\
&+\frac12\frac{\partial^2\tilde\rho}{\partial y^2}{\beta'_2}^2
-\frac{i\epsilon}{\hbar}v_{\mathrm R}(x)\tilde\rho
+\frac{i\epsilon}{\hbar}v_{\mathrm R}(y)\tilde\rho
-\frac{2M\gamma kT\epsilon}{\hbar}(x-y)^2\tilde\rho\Bigg].
\end{aligned}
\tag{5.8}\label{eq:5.8}
\end{equation}
式\eqref{eq:5.8}に現れる積分はすべて Fresnel 型であり、その主要な寄与が $\beta'_1=\beta'_2=0$ のまわりの微小な領域から来ることが分かっているので、積分範囲を $-\infty$ から $+\infty$ として評価できる。式\eqref{eq:5.8}の左辺を $\mathcal O(\epsilon)$ の項まで展開し、右辺の対応する項と等置すると、以下のことが分かる。

a) $\epsilon$ の零次の項から、規格化定数
\begin{equation}
A^2=\frac{2\pi\epsilon\hbar}{M}
\tag{5.9}\label{eq:5.9}
\end{equation}
が得られる。

b) $\epsilon$ の一次の項から、半古典的領域における $\tilde\rho$ の求める運動方程式
\begin{equation}
\begin{aligned}
\frac{\partial\tilde\rho}{\partial t}={}&
-\frac{\hbar}{2Mi}\frac{\partial^2\tilde\rho}{\partial x^2}
+\frac{\hbar}{2Mi}\frac{\partial^2\tilde\rho}{\partial y^2}
-\gamma(x-y)\frac{\partial\tilde\rho}{\partial x}
+\gamma(x-y)\frac{\partial\tilde\rho}{\partial y}
+\frac{v_{\mathrm R}(x)}{i\hbar}\tilde\rho\\
&-\frac{v_{\mathrm R}(y)}{i\hbar}\tilde\rho
-\frac{2M\gamma kT}{\hbar^2}(x-y)^2\tilde\rho
\end{aligned}
\tag{5.10}\label{eq:5.10}
\end{equation}
が得られる。

式\eqref{eq:5.10}が $\tilde\rho$ に対する最も一般的な方程式ではないことを、ここでもう一度強調しておきたい。この式が有効なのは $2kT\gtrsim\hbar\Omega\gg\hbar\omega_{\mathrm R}$ の場合に限られる。より一般的な結果を得ようとするなら、式\eqref{eq:3.38}の代わりに式\eqref{eq:5.1}を書くことは許されず、式\eqref{eq:5.10}の最後の項は時間積分を含むことになる（系の過去の履歴に依存することになる）。

以前に導かれたマスター方程式と比較するには、式\eqref{eq:5.10}の形は最も適したものではない。むしろ、表現に依存しない演算子方程式として書きたい。そのためには、
\begin{equation}
\begin{gathered}
\langle x\mid y\rangle=\delta(x-y),\qquad
\langle x\mid v(x')\mid y\rangle=v(y)\delta(x-y),\\
\langle x\mid p\mid y\rangle=-i\hbar\frac{\partial}{\partial x}\delta(x-y)
\end{gathered}
\tag{5.11}\label{eq:5.11}
\end{equation}
などの恒等式を用いればよい。これらを用いると、式\eqref{eq:5.10}は演算子方程式
% 原文 PDF p.22（誌面 p.608）。
\begin{equation}
\begin{aligned}
\frac{\partial\tilde\rho}{\partial t}={}&
\frac{1}{i\hbar}[H_{\mathrm R},\tilde\rho]
+\frac{\gamma}{2i\hbar}[\{p,x\},\tilde\rho]
-\frac{D}{\hbar^2}[x,[x,\tilde\rho]]\\
&+\frac{\Delta}{\hbar^2}\bigl([x,\tilde\rho p]-[p,\tilde\rho x]\bigr)
\end{aligned}
\tag{5.12}\label{eq:5.12}
\end{equation}
の座標表示であることを確かめられる。ここで $H_{\mathrm R}$ は系単独の繰り込まれた Hamiltonian、$D=\eta kT$、$\Delta=-i\hbar\gamma$ であり、$[\ ,\ ]$ は交換子、$\{\ ,\ \}$ は反交換子を表す。式\eqref{eq:5.12}は Dekker\cite{ref07} により、全く異なる方法で導かれている。実際には、彼の方程式は（$D$ 以外に）さらに二つの拡散項を含み、それらがある条件のもとで消えることを彼は示している。さらに、序論で述べたように、これらの追加項は位置の運動方程式に雑音源を含めたことに由来する。もう一つの相違は、Dekker の方程式が、$D$ にある有限の値を与えることで零温度においても有効なのに対し、我々の式は $T$ が高く $D=\eta kT$ の場合にしか有効でないことである。

本節の初めに、$\tilde\rho$ に対するマスター方程式が得られても、それを位相空間内の古典的分布の方程式と比較するには、さらにもう一段階が必要であると指摘した。前に述べたように、ここで必要なのは、
\begin{equation}
w(x,p,t)=\frac{1}{2\pi\hbar}\int_{-\infty}^{\infty}
\exp\!\left(\frac{ipy}{\hbar}\right)
\left\langle x-\frac y2\,\middle|\,\tilde\rho\,\middle|\,x+\frac y2\right\rangle\mathrm dy
\tag{5.13}\label{eq:5.13}
\end{equation}
で定義される Wigner 分布（すなわち $\tilde\rho$ の Wigner 変換）である。

Wigner 分布は系の純粋に量子的な記述を与えるものだが、半古典的領域に関心があるときに非常に重要になる。その理由は、Wigner の理論が系の古典的な位相空間内で直接記述されるためである。$\hbar\to0$ のとき、$w(x,p,t)$ は位相空間内の古典的分布に近づく。この主張の妥当性についての議論は、文献\cite{ref03,ref30}を参照されたい。

式\eqref{eq:5.12}の Wigner 変換を取り、恒等式\eqref{eq:5.11}を用いると、次式が得られる（詳細は文献\cite{ref07}を参照）。
\begin{equation}
\frac{\partial w}{\partial t}
=-\frac{\partial}{\partial x}(pw)
+\frac{\partial}{\partial p}(v'_{\mathrm R}w)
+2\gamma\frac{\partial}{\partial p}(pw)
+D\frac{\partial^2w}{\partial p^2}.
\tag{5.14}\label{eq:5.14}
\end{equation}
これは、系の縮約密度演算子の Wigner 変換の時間発展を記述する、よく知られた Fokker--Planck 方程式である。これは純粋に量子力学的な方程式である。しかし、$D=\eta kT$ としていることと整合させるためには、$w$ の式で $\hbar\to0$ の極限を取る必要がある。前に述べたように、この極限において $w(p,x,t)$ は古典的な位相空間分布に近づく。したがって、$\hbar\to0$ のとき、式\eqref{eq:5.14}は古典的な Brownian 粒子の位相空間分布の時間発展を記述すると結論できる。

% 原文 PDF p.23（誌面 p.609）。
また、式\eqref{eq:5.14}の右辺第三項は、式\eqref{eq:3.38}の指数の虚部に $\gamma$ に依存する項が存在することの直接の帰結であると分かる。したがって、この項が Brownian 粒子の古典的運動方程式に $\eta\dot x$ の形の力を出現させる原因である。

ここまでで得られたのは、伝播関数 $J$ に対する我々の量子力学的な式\eqref{eq:3.38}が、Brownian 粒子の古典的運動を記述すると考えられている方程式と完全に整合するということである。またこれは、式\eqref{eq:3.23}における $\rho_{\mathrm D}(\omega)C^2(\omega)$ の選択が非常に適切であり、現象論的な減衰定数 $\eta$ だけを用いて $J$ を記述できることも意味している。

\sectionlink{いくつかの応用}\label{sec:6}
本節では、いくつかの具体例、すなわち減衰単振動子と、その特殊な場合である減衰する自由粒子について、$\tilde\rho(x,y,t)$ の振る舞いを調べる。極端な量子極限（$T\to0$）にも関心があるので、ここでは高温という制限を外す。ただし、その前に、本論文で示した模型は何種類かの運動と両立することを指摘しておきたい。例えば、調和ポテンシャル $v(x)=\frac12M\omega^2x^2$ から出発するとしよう。この理論の微視的パラメータを適切に選ぶことにより、不足減衰運動と過減衰運動のいずれも記述できることが容易に示される。これは、座標結合の場合にも速度結合の場合にも成り立つ。両者の相違は、前者では $\gamma$ と比較すべき振動数が $\omega_{\mathrm R}^2\equiv\omega^2-(\Delta\omega)^2$ で与えられる繰り込まれた振動数 $\omega_{\mathrm R}$ であるのに対し、後者では $\omega_{\mathrm R}=\omega$ となることである。詳細は文献\cite{ref31}を参照されたい。

これから、調和振動子および自由粒子（$\omega_{\mathrm R}=0$）について式\eqref{eq:3.38}を評価する。まず、式\eqref{eq:3.38}を
\begin{equation}
J(x_{\mathrm f},y_{\mathrm f},t;x_{\mathrm i},y_{\mathrm i},0)
=\iint\mathcal Dx\,\mathcal Dy\,
\exp\!\left(\frac{i}{\hbar}\tilde S[x,y]\right)
\exp\!\left(-\frac1\hbar\phi[x,y]\right)
\tag{6.1}\label{eq:6.1}
\end{equation}
と書き直す。ここで
\begin{equation}
\tilde S[x,y]=\int_0^tL(x,\dot x,y,\dot y)\,\mathrm d\tau
-\int_0^tM\gamma x\dot x\,\mathrm d\tau
+\int_0^tM\gamma y\dot y\,\mathrm d\tau,
\tag{6.2}\label{eq:6.2}
\end{equation}
\begin{equation}
L=\frac12M\dot x^2-\frac12M\dot y^2
-\frac12M\omega_{\mathrm R}^2x^2+\frac12M\omega_{\mathrm R}^2y^2
-M\gamma x\dot y+M\gamma y\dot x
\tag{6.3}\label{eq:6.3}
\end{equation}
および
\begin{equation}
\begin{aligned}
\phi[x,y]={}&\frac{2M\gamma}{\pi}\int_0^\Omega
\nu\coth\!\left(\frac{\hbar\nu}{2kT}\right)
\int_0^t\int_0^\tau[x(\tau)-y(\tau)]\\
&\qquad\times\cos\nu(\tau-s)[x(s)-y(s)]\,\mathrm ds\,\mathrm d\tau\,\mathrm d\nu
\end{aligned}
\tag{6.4}\label{eq:6.4}
\end{equation}
である。

% 原文 PDF p.24（誌面 p.610）。
式\eqref{eq:6.1}を評価する手順は簡単である。$\tilde S[x,y]$ を極値にする二次元経路のまわりで被積分関数全体を展開する。この経路は
\begin{equation}
\frac{\delta\tilde S}{\delta x}
=\frac{\mathrm d}{\mathrm d\tau}\frac{\partial L}{\partial\dot x}
-\frac{\partial L}{\partial x}
=M\ddot x+M\gamma\dot y+M\omega_{\mathrm R}^2x=0,
\tag{6.5}\label{eq:6.5}
\end{equation}
\begin{equation}
\frac{\delta\tilde S}{\delta y}
=\frac{\mathrm d}{\mathrm d\tau}\frac{\partial L}{\partial\dot y}
-\frac{\partial L}{\partial y}
=M\ddot y+M\gamma\dot x+M\omega_{\mathrm R}^2y=0
\tag{6.6}\label{eq:6.6}
\end{equation}
から求められる。式\eqref{eq:6.2}の最後の二項は端点の変分だけを含むためである。式\eqref{eq:6.5}と\eqref{eq:6.6}を解くには、新しい変数を導入すると便利である。そこで
\begin{equation}
X(\tau)=x(\tau)+y(\tau),\qquad
\xi(\tau)=x(\tau)-y(\tau)
\tag{6.7}\label{eq:6.7}
\end{equation}
と定義すると、方程式は
\begin{equation}
\ddot X+2\gamma\dot X+\omega_{\mathrm R}^2X=0,
\tag{6.8}\label{eq:6.8}
\end{equation}
\begin{equation}
\ddot\xi-2\gamma\dot\xi+\omega_{\mathrm R}^2\xi=0
\tag{6.9}\label{eq:6.9}
\end{equation}
となる。

$X(t)=X_{\mathrm f}$、$\xi(t)=\xi_{\mathrm f}$、$X(0)=X_{\mathrm i}$、$\xi(0)=\xi_{\mathrm i}$ のとき、不足減衰調和振動子（以下 UDHO）について、その解は
\begin{equation}
\tilde X(\tau)=(\sin\omega t)^{-1}
\bigl\{X_{\mathrm f}e^{\gamma t}\sin\omega\tau
+X_{\mathrm i}\sin\omega(t-\tau)\bigr\}e^{-\gamma\tau},
\tag{6.10}\label{eq:6.10}
\end{equation}
\begin{equation}
\tilde\xi(\tau)=(\sin\omega t)^{-1}
\bigl\{\xi_{\mathrm f}e^{-\gamma t}\sin\omega\tau
+\xi_{\mathrm i}\sin\omega(t-\tau)\bigr\}e^{\gamma\tau}
\tag{6.11}\label{eq:6.11}
\end{equation}
であり、ここで $\omega^2\equiv\omega_{\mathrm R}^2-\gamma^2$ である。過減衰の場合（ODHO）の解は、式\eqref{eq:6.10}と\eqref{eq:6.11}で $\omega$ を $i\tilde\omega$ に置き換えることにより得られる。ここで $\tilde\omega^2\equiv\gamma^2-\omega_{\mathrm R}^2$ である。自由粒子（FP）については $\omega$ を $i\gamma$ に置き換える。実際、これらの置換は以後の式にも常に適用できるが、その一部（FP の式）では、必要に応じて $\omega_{\mathrm R}\to0$ の極限を慎重に取る必要があることに注意されたい。このことを念頭に置き、本節の式は UDHO に対してだけ記すことにする。

経路\eqref{eq:6.10}--\eqref{eq:6.11}に沿って計算した作用 $\tilde S[X,\xi]$ は、
\begin{equation}
\begin{aligned}
\tilde S_{\mathrm{cl}}={}&K(t)[X_{\mathrm f}\xi_{\mathrm f}+X_{\mathrm i}\xi_{\mathrm i}]
-L(t)X_{\mathrm i}\xi_{\mathrm f}-N(t)X_{\mathrm f}\xi_{\mathrm i}\\
&-\frac{M\gamma}{2}[X_{\mathrm f}\xi_{\mathrm f}-X_{\mathrm i}\xi_{\mathrm i}]
\end{aligned}
\tag{6.12}\label{eq:6.12}
\end{equation}
と表せることが分かる。ここで
\begin{equation}
K(t)=\frac{M\omega}{2}\cot\omega t,\qquad
L(t)=\frac{M\omega e^{-\gamma t}}{2\sin\omega t},\qquad
N(t)=\frac{M\omega e^{\gamma t}}{2\sin\omega t}
\tag{6.13}\label{eq:6.13}
\end{equation}
である。

ここで、平行移動した経路
\begin{equation}
\xi'(\tau)=\xi(\tau)-\tilde\xi(\tau),\qquad
X'(\tau)=X(\tau)-\tilde X(\tau)
\tag{6.14}\label{eq:6.14}
\end{equation}
を用いて汎関数積分を書こう。
% 原文 PDF p.25（誌面 p.611）。
境界条件は $\xi'(0)=\xi'(t)=0$ および $X'(0)=X'(t)=0$ である。すると、
\begin{equation}
\begin{aligned}
J(X_{\mathrm i},\xi_{\mathrm f},t;X_{\mathrm i},\xi_{\mathrm i},0)
={}&\exp\!\left(\frac{i}{\hbar}\tilde S_{\mathrm{cl}}\right)
\exp\!\left[-\frac1\hbar\{A(t)\xi_{\mathrm f}^2
+B(t)\xi_{\mathrm f}\xi_{\mathrm i}+C(t)\xi_{\mathrm i}^2\}\right]\\
&\times G(X_{\mathrm f},\xi_{\mathrm f},t;X_{\mathrm i},\xi_{\mathrm i},0)
\end{aligned}
\tag{6.15}\label{eq:6.15}
\end{equation}
が得られる。ここで
\begin{equation}
A(t)=\frac{M\gamma}{\pi}\int_0^\Omega
\mathrm d\nu\,\nu\coth\!\left(\frac{\hbar\nu}{2kT}\right)A_\nu(t),
\tag{6.16}\label{eq:6.16}
\end{equation}
\begin{equation}
A_\nu(t)=\frac{e^{-2\gamma t}}{\sin^2\omega t}
\int_0^t\int_0^t\sin\omega\tau\cos\nu(\tau-s)\sin\omega s\,
e^{\gamma(\tau+s)}\,\mathrm d\tau\,\mathrm ds,
\tag{6.17}\label{eq:6.17}
\end{equation}
\begin{equation}
B(t)=\frac{M\gamma}{\pi}\int_0^\Omega
\mathrm d\nu\,\nu\coth\!\left(\frac{\hbar\nu}{2kT}\right)B_\nu(t),
\tag{6.18}\label{eq:6.18}
\end{equation}
\begin{equation}
B_\nu(t)=\frac{2e^{-\gamma t}}{\sin^2\omega t}
\int_0^t\int_0^t\sin\omega\tau\cos\nu(\tau-s)\sin\omega(t-s)\,
e^{\gamma(\tau+s)}\,\mathrm d\tau\,\mathrm ds,
\tag{6.19}\label{eq:6.19}
\end{equation}
\begin{equation}
C(t)=\frac{M\gamma}{\pi}\int_0^\Omega
\mathrm d\nu\,\nu\coth\!\left(\frac{\hbar\nu}{2kT}\right)C_\nu(t),
\tag{6.20}\label{eq:6.20}
\end{equation}
\begin{equation}
C_\nu(t)=\frac1{\sin^2\omega t}
\int_0^t\int_0^t\sin\omega(t-\tau)\cos\nu(\tau-s)\sin\omega(t-s)\,
e^{\gamma(\tau+s)}\,\mathrm d\tau\,\mathrm ds,
\tag{6.21}\label{eq:6.21}
\end{equation}
\begin{equation}
\begin{aligned}
G(X_{\mathrm f},\xi_{\mathrm f},t;X_{\mathrm i},\xi_{\mathrm i},0)
={}&\iint\frac{\mathcal D\xi'\,\mathcal DX'}2
\exp\!\left(\frac{i}{\hbar}\tilde S[X',\xi']\right)
\exp\!\left(-\frac2\hbar\phi_T[\xi',\xi']\right)\\
&\times\exp\!\left(\frac1\hbar\phi_T[\tilde\xi,\xi']\right)
\end{aligned}
\tag{6.22}\label{eq:6.22}
\end{equation}
および
\begin{equation}
\phi_T[\tilde\xi,\xi']=\frac{M\gamma}{\pi}\int_0^\Omega
\mathrm d\nu\,\nu\coth\!\left(\frac{\hbar\nu}{2kT}\right)
\int_0^t\int_0^t\tilde\xi(\tau)\cos\nu(\tau-s)\xi'(s)\,\mathrm d\tau\,\mathrm ds
\tag{6.23}\label{eq:6.23}
\end{equation}
である。

残っているのは汎関数積分\eqref{eq:6.22}の評価である。$\xi'(\tau)$ と $X'(\tau)$ はどちらも $\tau=0$ と $\tau=t$ で消えるので、同じ境界条件を満たす直交系を用いて展開できる\cite{ref22}。
% 原文 PDF p.26（誌面 p.612）。
\begin{equation}
X'(\tau)=\sum_nX_n\sin\omega_n\tau,\qquad
\xi'(\tau)=\sum_m\xi_m\sin\omega_m\tau,
\tag{6.24}\label{eq:6.24}
\end{equation}
ここで $\omega_n=n\pi/t$ である。式\eqref{eq:6.24}を式\eqref{eq:6.22}に代入し、いくつかの簡単な積分を行うと（同様の手順は文献\cite{ref22}を参照）、
\begin{equation}
G(X_{\mathrm f},\xi_{\mathrm f},t;X_{\mathrm i},\xi_{\mathrm i},0)
=F^2(t)=\lim_{N\to\infty}\frac{4\tilde C}{Mt}
(2\pi\hbar)^N\prod_{i=1}^N\frac1{\omega_i^2-\omega_{\mathrm R}^2}
\tag{6.25}\label{eq:6.25}
\end{equation}
が得られる。ここで $\tilde C$ は規格化定数である。この結果は、減衰がない場合の波動関数の振幅の二乗と正確に等しい。したがって、減衰が変えるのは $J(X_{\mathrm f},\xi_{\mathrm f},t;X_{\mathrm i},\xi_{\mathrm i},0)$ の指数だけであり、その最終的な形は
\begin{equation}
\begin{aligned}
J={}&F^2(t)\exp\!\Bigg[\frac{i}{\hbar}\Bigg\{
\left[K(t)-\frac{M\gamma}{2}\right]X_{\mathrm f}\xi_{\mathrm f}
+\left[K(t)+\frac{M\gamma}{2}\right]X_{\mathrm i}\xi_{\mathrm i}\\
&\hspace{38mm}-L(t)X_{\mathrm i}\xi_{\mathrm f}
-N(t)X_{\mathrm f}\xi_{\mathrm i}\Bigg\}\Bigg]
\exp\!\left[-\frac1\hbar\{A(t)\xi_{\mathrm f}^2
+B(t)\xi_{\mathrm f}\xi_{\mathrm i}+C(t)\xi_{\mathrm i}^2\}\right]
\end{aligned}
\tag{6.26}\label{eq:6.26}
\end{equation}
となる。

これで、関心のある系の縮約密度演算子の時間発展を計算するために必要な道具が得られた。例えば、初めに系が、原点を中心とし初期運動量 $p$、幅 $\sigma$ の波束で記述される純粋状態にあるとしよう。このとき縮約密度演算子は
\begin{equation}
\tilde\rho(X_{\mathrm i},\xi_{\mathrm i},0)
=(2\pi\sigma^2)^{-1/2}
\exp\!\left(\frac{ip\xi_{\mathrm i}}{\hbar}\right)
\exp\!\left(-\frac{X_{\mathrm i}^2+\xi_{\mathrm i}^2}{8\sigma^2}\right)
\tag{6.27}\label{eq:6.27}
\end{equation}
であり、その時間発展は
\begin{equation}
\tilde\rho(X_{\mathrm f},\xi_{\mathrm f},t)
=\iint\mathrm dX_{\mathrm i}\,\mathrm d\xi_{\mathrm i}\,
J(X_{\mathrm f},\xi_{\mathrm f},t;X_{\mathrm i},\xi_{\mathrm i},0)
\tilde\rho(X_{\mathrm i},\xi_{\mathrm i},0)
\tag{6.28}\label{eq:6.28}
\end{equation}
と書ける。これは式\eqref{eq:6.27}と\eqref{eq:6.26}を代入することで評価できる。式\eqref{eq:6.28}の積分はすべて Gaussian 積分であり、非常に簡単である。しかし、その簡単さにもかかわらず、中間段階で得られる式はかなり長いため、ここでは最終結果だけを示す。
\begin{equation}
\begin{aligned}
\tilde\rho(X_{\mathrm f},\xi_{\mathrm f},t)
={}&F^2(t)\left[\frac{\pi^2\hbar^2}{2\sigma^2K_1+\hbar C_1}\right]
\exp\!\left[-\frac{N^2}{2\sigma^2K_1^2+4\hbar C_1^2}
\left(X_{\mathrm f}-\frac p{2N}\right)^2\right]\\
&\times\exp\!\left[-\left\{\frac A\hbar+\frac{2\sigma^2L^2}{\hbar^2}
-\frac{(4\sigma^2K_1L-B\hbar)^2}{\hbar^2(8\sigma^2K_1^2+4\hbar C_1)}\right\}\xi_{\mathrm f}^2\right]\\
&\times\exp\!\left[\frac1\hbar\left\{K_2X_{\mathrm f}\xi_{\mathrm f}
-\frac{(4\sigma^2K_1L-B\hbar)N}{4\sigma^2K_1^2+2\hbar C_1}
\left(X_{\mathrm f}-\frac pN\right)\xi_{\mathrm f}\right\}\right].
\end{aligned}
\tag{6.29}\label{eq:6.29}
\end{equation}
% 原文 PDF p.27（誌面 p.613）。
ここで、$A$、$B$、$C_1$、$K_1$、$K_2$、$L$、$N$ はすべて時間の関数であることに注意し、
\begin{equation}
C_1(t)\equiv C(t)+\frac{\hbar}{8\sigma^2},\qquad
K_1(t)\equiv K(t)+\frac{M\gamma}{2},\qquad
K_2(t)\equiv K(t)-\frac{M\gamma}{2}
\tag{6.30}\label{eq:6.30}
\end{equation}
と定義した。

特に興味深いのは、式\eqref{eq:6.29}において $\xi_{\mathrm f}=x_{\mathrm f}-y_{\mathrm f}=0$、すなわち $X_{\mathrm f}=x_{\mathrm f}+y_{\mathrm f}=2x_{\mathrm f}$ とすることで得られる式
\begin{equation}
\tilde\rho(x_{\mathrm f},x_{\mathrm f},t)
=F^2(t)\left[\frac{\pi^2\hbar^2}{2\sigma^2K_1^2+\hbar C_1}\right]^{1/2}
\exp\!\left[-\frac{N^2}{2\sigma^2K_1^2+\hbar C_1}
\left(x_{\mathrm f}-\frac p{2N}\right)^2\right]
\tag{6.31}\label{eq:6.31}
\end{equation}
である。これは、初期の $|\psi(x)|^2$ が時間とともにどのように発展するかを示している。その中心は
\begin{equation}
x_0(t)=\frac p{2N(t)}
\tag{6.32}\label{eq:6.32}
\end{equation}
で与えられる経路に従う。これは減衰する古典粒子の経路であり、例えば UDHO では
\begin{equation}
x_0(t)=\frac p{M\omega}\sin\omega t\,e^{-\gamma t}
\tag{6.33}\label{eq:6.33}
\end{equation}
となる。一方、その幅は
\begin{equation}
\sigma^2(t)=\frac12\left(\frac{2\sigma^2K_1^2(t)+\hbar C_1(t)}{N^2(t)}\right)
\tag{6.34}\label{eq:6.34}
\end{equation}
に従う。

この式は、任意の時刻 $t$ における $\tilde\rho(x,x,t)$ の幅を与える。これを無限時間で調べると、興味深い結果が得られる。手順は単純だが、非常に煩雑である。式\eqref{eq:6.13}、\eqref{eq:6.20}、\eqref{eq:6.21}、\eqref{eq:6.30}を式\eqref{eq:6.34}に代入すると、
\begin{equation}
\begin{aligned}
\sigma^2(t\to\infty)
&=\left\langle\bigl(x-x_0(\infty)\bigr)^2\right\rangle\\
&=\frac{\hbar}{\pi}\int_0^\infty\mathrm d\nu\,
\coth\!\left(\frac{\hbar\nu}{2kT}\right)
\left(\frac1M\frac{2\gamma\nu}{(\omega_{\mathrm R}^2-\nu^2)^2+4\gamma^2\nu^2}\right)
\end{aligned}
\tag{6.35}\label{eq:6.35}
\end{equation}
が得られる。ここでは積分の上限で $\Omega\to\infty$ の極限を取った。式\eqref{eq:6.35}は非常によく知られたものである。振動数 $\omega_{\mathrm R}$ の調和振動子の外力 $F(t)$ に対する応答の虚部は
\begin{equation}
\chi''(\nu)=\frac1M\frac{2\gamma\nu}{(\nu^2-\omega_{\mathrm R}^2)^2+4\gamma^2\nu^2}
\tag{6.36}\label{eq:6.36}
\end{equation}
で与えられる。したがって、
\begin{equation}
\sigma^2(\infty)=\frac{\hbar}{\pi}\int_0^\infty\mathrm d\nu\,
\coth\!\left(\frac{\hbar\nu}{2kT}\right)\chi''(\nu)
\tag{6.37}\label{eq:6.37}
\end{equation}
となり、これはよく知られた揺動散逸定理である\cite{ref32}。

% 原文 PDF p.28（誌面 p.614）。
\sectionlink{結論}\label{sec:7}
熱浴と相互作用する粒子の具体的な模型を用いて、その粒子の縮約密度演算子の伝播関数に対する式\eqref{eq:3.38}を導いた。前節で示したように、ある状況では、任意の時刻 $t$、あらゆる温度における $\tilde\rho(x,y,t)$ の閉じた式を得ることができた。波束 $\tilde\rho(x,x,t)$ の中心は古典的な Brownian 粒子の運動に従い（式\eqref{eq:6.33}を参照）、その無限時間極限における幅は揺動散逸定理\eqref{eq:6.35}と一致することを示した。後者の結果について、もう少し詳しく考察できる。

まず、式\eqref{eq:6.35}は文献\cite{ref27}と一致しない。そこでは、自由粒子の波束の幅は無限時間で「凍結する」と主張されている。式\eqref{eq:6.35}で $\omega_{\mathrm R}\to0$ とすると、$T\to0$ のとき、式\eqref{eq:6.35}は $\nu\to0$ において $(-\ln\nu)$ のように発散することを示せる。これは、$t\to\infty$ のとき $\sigma^2(t)$ が $\ln t$ のように振る舞うということに等しい。したがって、$\sigma^2(t)$ の広がりは減衰がない場合（$\sigma^2(t)=\sqrt{A+Bt^2}$）より遅いが、凍結することはない。

我々は、両者の方法のこの相違は、我々と異なり文献\cite{ref27}の著者たちがいかなる拡散効果も考慮していないことに起因すると考えている。式\eqref{eq:6.1}で $\phi[x,y]=0$ とすることにより、二つの方法の物理的解釈が全く異なるにもかかわらず、彼らの結果を再現できる。

Morse と Feshbach の Lagrangian から出発して Brownian 運動を量子化しようとする他の試みにおいても、拡散項は省略されていた（例えば、文献\cite{ref33,ref34}、あるいは文献\cite{ref35}の第3節と第6節を参照）。

第二の比較は、減衰が極めて小さい極限（$\gamma\to0$）を調べるときに現れる。この場合、$t\to\infty$ における最終的な幅は、任意の温度で
\begin{equation}
\sigma^2(\infty)=\frac{\hbar}{2M\omega_{\mathrm R}}
\coth\!\left(\frac{\hbar\omega_{\mathrm R}}{2kT}\right)
\tag{7.1}\label{eq:7.1}
\end{equation}
と求められる。これは、平衡量子統計力学のよく知られた結果である\cite{ref22}。式\eqref{eq:3.38}で与えられる厳密な伝播関数 $J$ の代わりに、式\eqref{eq:5.1}において $2MkT\gamma$ を $\hbar M\gamma\omega_{\mathrm R}\coth(\hbar\omega_{\mathrm R}/2kT)$ に置き換えたものを用いても、同じ式が得られるはずである。言い換えると、（運動量空間における）拡散定数は
\begin{equation}
D=M\gamma\omega_{\mathrm R}\hbar
\coth\!\left(\frac{\hbar\omega_{\mathrm R}}{2kT}\right)
\tag{7.2}\label{eq:7.2}
\end{equation}
で与えられ、あらゆる温度で有効であるはずだ。$T\to\infty$ のとき、式\eqref{eq:7.2}は古典的な関係 $D=2M\gamma kT$ を与え、一方、$T\to0$ のとき拡散定数は
\begin{equation}
D=M\gamma\hbar\omega_{\mathrm R}=\frac12\eta\hbar\omega_{\mathrm R}
\tag{7.3}\label{eq:7.3}
\end{equation}
となる。これは Dekker\cite{ref07} が見いだした結果である。

% 原文 PDF p.29（誌面 p.615）。
序論で指摘したように、Dekker は位置演算子と運動量演算子の運動方程式に明示的な雑音源を用い、複素変数に正準量子化を適用した（ただし、式\eqref{eq:7.3}の別の導出については文献\cite{ref36}を参照）。確かに、こうした操作を許す微視的模型は存在する。しかし、非常に弱い減衰の極限では、これらの方程式の解は通常の方程式（運動量の方程式だけに雑音源があるもの）の解と同じになることを示せる\cite{ref37}。

式\eqref{eq:7.2}で提案した工夫は、有限の $\gamma$ に対しては正しくないように思われる。この工夫は、式\eqref{eq:3.9}の影響汎関数の指数に現れる核 $\alpha_{\mathrm R}(t-t')$ を、常に温度の関数と $(t-t')$ のデルタ関数との積で置き換えられることを意味する。しかし、我々の具体的な模型ではそうならない。式\eqref{eq:7.2}が任意の $\gamma$ に対して正しくなるような別の熱浴模型を見いだせるとしても、我々の模型は、この手順が一般的でないことを示す反例となる。

式\eqref{eq:7.2}の形の拡散定数は Papadopoulos\cite{ref28} と Svin'in\cite{ref38} によって提案されている。しかし、ここでも、彼らの理論が有効なのは弱い減衰の極限だけである。これは、その式が成り立つための必要条件だからである。

実際、量子的領域においてさえも、Markoff 型の方程式（時間核を含まないもの）で平衡への緩和を記述しようとすることは、多くの物理学者にとって長年の夢であった。しかし、我々はその可能性にやや懐疑的である。我々が示したように、簡単な模型でも、全く Markoff 型でないまま、平衡統計力学と線形応答理論の両方を漸近的に再現できるためである。Ford ら\cite{ref16}も、量子極限における Markoff 型という仮定について同じ結論に達している。ここで論じているのは拡散項だけであることに注意されたい。我々のドリフト項（$\gamma$ を含む項）は常に Markoff 型である。後者に対する非 Markoff 型の補正を求めるには、Brownian 粒子の短時間の振る舞いを記述するのに重要な総和則による $\chi''(\nu)$ の修正を調べる必要がある。

最後に、系と環境の結合について、一見非常に恣意的に見える模型を採用したが、結果の主要な特徴は、その詳細にはあまり左右されないはずであることを強調しておきたい。特に、量 $\rho_{\mathrm D}(\omega)C^2(\omega)$ に対する式\eqref{eq:3.23}の選択を変更しても、系の特徴的な振動数と同程度の $\omega$ に対して式\eqref{eq:3.23}の形が引き続き成り立ちさえすれば、結果に定性的な違いはほとんど生じないはずである。実際に影響を受けるのは振動数のずれ $(\Delta\omega)^2$ だけであり、いずれにせよ、これは直接には観測できない（トンネル問題における類似の点についての文献\cite{ref21}の議論を参照）。

\sectionlink*{謝辞}
有益な議論、書簡のやり取り、およびプレプリントを提供してくださった G. Barton 博士、P.L. Knight 博士、A.B. Pippard 教授、K.L. Sebastian 博士、D.J. Wallace 教授、A.F.R. de Toledo Piza 教授に感謝する。

\clearpage
\renewcommand{\refname}{\hyperlink{toc}{参考文献}}
\bibliographystyle{utphys}
\bibliography{ref}
\end{document}
